\documentclass[11pt]{article}
\usepackage[margin=1in]{geometry}
\usepackage[T1]{fontenc}\usepackage[utf8]{inputenc}\usepackage{lmodern}
\usepackage{amsmath,amssymb,amsthm,mathtools,microtype}
\usepackage[colorlinks=true,linkcolor=blue,urlcolor=blue]{hyperref}
\hypersetup{pdftitle={A depth obstruction for Gaussian Bayesian tanh networks from a maximum-amplitude function prior},pdfauthor={ChatGPT-6 Astra; Claude Opus 5.5; Kaarel H\"anni; Dmitry Vaintrob}}
\newtheorem{theorem}{Theorem}[section]\newtheorem{lemma}[theorem]{Lemma}\newtheorem{corollary}[theorem]{Corollary}\newtheorem{proposition}[theorem]{Proposition}
\newcommand{\E}{\mathbb E}\newcommand{\R}{\mathbb R}\newcommand{\KL}{\operatorname{KL}}
\title{A depth obstruction for Gaussian Bayesian tanh networks\\from a maximum-amplitude function prior}
\author{ChatGPT-6 Astra \and Claude Opus 5.5 \and Kaarel H\"anni \and Dmitry Vaintrob\thanks{Authors are listed in alphabetical order.}}
\date{17 November 2026}
\begin{document}\maketitle
\begin{abstract}
We prove a fixed-accuracy learning obstruction for deep Bayesian tanh networks across every Gaussian scale at or stronger than mean field. For $H$ hidden layers, one bounded noisy regression problem on a unit sphere in dimension $O(\sqrt H)$ has a width-two, one-hidden-layer conditional mean, yet after $n=\exp(\Omega(\sqrt H))$ IID samples every width $N\ge n^C$ has fixed positive excess squared risk. One sample event covers all prior scales $0<s\le1/N$, all Gaussian-likelihood variances, clipped prediction conventions, and arbitrary data-dependent mixtures of the posterior means. The independent weight variances are $Ns/d$ in the first layer, $s$ internally, and $s/N$ in the readout. For the family with $Ns\le1-\Delta$, the sample obstruction strengthens to $\exp(\Omega(\sqrt H+\Delta H))$ in dimension $O(\sqrt H+\Delta H)$, uniformly over $0\le\Delta<1$. The analytic mechanism is a function prior that prices small finite-domain functions mainly by maximum amplitude. At canonical mean field its precision-normalized relative-entropy cost is nearly $4\|f\|_\infty$, with a matching lower bound; smaller variance prefactors add a depth baseline that cancels in the statistical comparison. The upper construction uses rank-one conditional covariance changes, while the learner's prior remains unchanged. The obstruction also holds with correctly specified Gaussian observation noise and arbitrary independent Gaussian affine output components. Final clipping or tanh inside the likelihood is covered separately. The same depth-gap crossover holds for each extension with its stated prediction convention. The results are finite-width statements and use no prior over network sizes.
\end{abstract}

\section{The Gaussian family and the main learning obstruction}
A shallow conditional mean can be easy to represent while remaining hard for a much deeper Bayesian learner to recover from noisy data. The mechanism here is that sufficient depth makes arbitrary small-amplitude patterns on a large finite input set nearly as cheap under the prior as the signal. Low-temperature posteriors then retain noisy empirical signs, while high-temperature posteriors shrink toward zero; both behaviors give a common excess-risk obstruction.

A scalar readout follows $H\ge1$ hidden tanh layers of common width $N$:
\[
 h^1(x)=\tanh(W^1(1,x)),\qquad
 h^\ell(x)=\tanh(W^\ell h^{\ell-1}(x)),\qquad
 z_w(x)=v^{\mathsf T}h^H(x),\qquad x\in\R^d.
\]
There are no internal biases. The Gaussian family has independent centered coordinates of variances $Ns/d$, $s$, and $s/N$ in the first, internal, and readout layers, respectively, with $0<s\le1/N$. Canonical mean field is $s=1/N$. For the analytic development until Section~\ref{sec:joint-prefactor-entropy}, write $\mu$ for that canonical prior, whose variances are
\begin{equation}
 \operatorname{Var}(W^1_{ji})=1/d,\qquad
 \operatorname{Var}(W^\ell_{ji})=1/N\quad(\ell\ge2),\qquad
 \operatorname{Var}(v_j)=1/N^2.
 \label{eq:canonical-prior}
\end{equation}
The prior has full support and does not depend on a target function. A comparison law $\rho\ll\mu$ is used only to upper-bound variational objectives or prior-neighborhood costs. It is not a replacement prior.

For an integer $D\ge1$, set $d=D+2$ and
\begin{equation}
 \mathcal U_D=\left\{\left(T,\frac z{2\sqrt D},\sqrt{3/4-T^2}\right):
 T\in\{\pm1/4,\pm3/4\},\ z\in\{-1,1\}^D\right\}
 \subset\mathbb S^{d-1}.
 \label{eq:unit-input-set}
\end{equation}
This set has $m=4\cdot2^D$ points. Its first coordinate takes four fixed values; the remaining Boolean coordinates index many independently prescribed outputs. We write $\|f\|_\infty=\max_{x\in\mathcal U_D}|f(x)|$.
For a prediction $f$ and a sample $S=((X_i,Y_i))_{i=1}^n$, define the population and empirical squared losses by
\[
 R(f)=\E[(f(X)-Y)^2],\qquad
 \widehat R_S(f)=\frac1n\sum_{i=1}^n(f(X_i)-Y_i)^2.
\]


\begin{theorem}[A common obstruction at every mean-field-or-stronger scale]
\label{thm:all-scale-learning}
There are absolute constants $a,c,C>0$ and an integer $H_0$ such that the following holds for every $H\ge H_0$. Put
\[
 D=\lfloor c\sqrt H\rfloor,\qquad d=D+2,\qquad n=4\cdot2^D,
\]
and let $X$ be uniform on the unit-sphere set $\mathcal U_D$. Let
\[
 Y=g(X)+a\epsilon,\qquad
 g(X)=\frac a4(2+\tanh X_1),\qquad
 \Pr(\epsilon=1)=\Pr(\epsilon=-1)=\frac12,
\]
with $\epsilon$ independent of $X$. Consider a width-$N$, depth-$H$ tanh network with a linear readout, independent centered Gaussian weights of variances
\[
 \frac{Ns}{d},\qquad s,\qquad\frac sN
\]
in its first, internal, and readout layers, respectively, and Gaussian likelihood
\[
 \Pi(dw)\propto\exp\{-\widehat R_S(z_w)/(2\lambda s)\}\,\mu_{N,s}(dw).
\]
The corresponding observation variance is $n\lambda s$. With probability at least $9/10$ over $n$ IID observations, simultaneously for every integer $N\ge n^C$, every $0<s\le1/N$, and every $\lambda>0$, the ordinary posterior mean has squared population loss at least
\[
 R(\widehat f)\ge a^2+\frac{a^2}{4096}.
\]
The same statement holds for the mean of posterior draws clipped to $[-1,1]$, clipping the ordinary posterior mean, and arbitrary well-defined data-dependent probability mixtures over these predictions, widths, scales, and temperatures, including clipping after mixing.

The conditional mean has width two and one hidden layer, with bounded coefficients independent of $H$. The sample count is $\exp(\Omega(\sqrt H))$, and the width threshold is a fixed polynomial in that count. The prior scale may depend arbitrarily on depth, width, and data inside the stated interval.
\end{theorem}

The term $a^2$ is the irreducible loss from the independent label noise. The theorem therefore gives a fixed positive excess loss. In particular it includes all canonical exponents $s=N^{-\gamma}$, $\gamma\ge1$, on the same event. It also includes strong-prior sequences at every width above the stated threshold. It concerns the Gaussian regression likelihood displayed in the theorem; the clipping operations in its statement are prediction conventions applied after forming that posterior. Theorem~\ref{thm:all-scale-affine-likelihood} allows an arbitrary independent Gaussian affine output component; Corollary~\ref{cor:all-scale-gaussian-affine} also replaces the bounded noise by correctly specified Gaussian noise. Corollary~\ref{cor:all-scale-bounded-likelihood} treats a final clip or tanh inside the likelihood. These extensions retain the same sample and width orders, with different absolute excess-risk constants. Theorem~\ref{thm:prefactor-window} additionally quantifies the transition below critical mean field: restricting $Ns\le1-\Delta$ strengthens the sample obstruction to $\exp(\Omega(\sqrt H+\Delta H))$, with constants uniform in the possibly depth-dependent gap $\Delta$. Corollary~\ref{cor:prefactor-window-extensions} verifies the same crossover for each observation-model extension with its stated prediction scope.


\subsection{The critical function-prior estimate}
\begin{theorem}[Nearly maximum-amplitude prior cost]
\label{thm:canonical-function-memorizer}
There are absolute $A_0,C_0,H_0>0$ such that the following holds. Let $H\ge H_0$, $1\le D\le\sqrt H/2$, and $0<A\le A_0$. For every tolerance $\varepsilon>0$, every width
\[
 N\ge C_0(1+\varepsilon^{-1})^{C_0}H^{C_0}e^{C_0D},
\]
and every function $f:\mathcal U_D\to[-A,A]$ admit a law $\rho\ll\mu$ with
\begin{align}
 \frac2N\KL(\rho\|\mu)
 &\le4\|f\|_\infty+C_0A^2+
 C_0\frac{(D+1)^2}{H}+\frac{2\log2}{N},
 \label{eq:canonical-function-cost}\\
 \E_\rho\max_{x\in\mathcal U_D}|z_w(x)-f(x)|^2&\le\varepsilon^2.
 \label{eq:canonical-function-approx}
\end{align}
Conversely, for every $0<\delta<1$ there is $C_\delta<\infty$ such that any comparison satisfying \eqref{eq:canonical-function-approx} obeys
\begin{equation}
 \frac2N\KL(\rho\|\mu)
 \ge(4-\delta)(\|f\|_\infty-\varepsilon)
       -\frac{C_\delta}{H}-\frac{2\log(2m)}N.
 \label{eq:canonical-function-lower}
\end{equation}
\end{theorem}

For every fixed $\delta>0$, choose a sufficiently small fixed $A>0$, then a sufficiently small fixed $c>0$. At $D=\lfloor c\sqrt H\rfloor$, the upper cost is at most $4\|f\|_\infty+\delta A$ for all large enough $H$. Thus the same leading cost covers every function on $\exp(\Omega(\sqrt H))$ inputs. The theorem concerns this depth and width regime; it is not an assertion that circuit complexity never influences a neural function prior.

The function-prior estimate first gives the canonical-scale learning argument in Theorem~\ref{thm:canonical-learning}. A joint depth-and-amplitude entropy bound then extends the argument to every smaller variance prefactor, proving the main theorem above. Thus a sample guarantee polynomial jointly in learner depth, input dimension, and teacher complexity, uniformly over learner width, is impossible throughout this unrestricted Gaussian family. The result does not exclude bounds depending polynomially on learner width, and fixed-depth guarantees remain possible.

The proof has two parts. Exponential moments of the native hidden-variance chain give the lower bound. For the upper bound, a Gaussian feature kernel first encodes any finite function in a neuron-space direction. Covariance spikes amplify that direction with total cost $O((D+1)^2/H)$, and a final variance-recovery phase costs nearly four times the requested amplitude. We first establish the scalar estimates and a simpler one-input construction, then give the finite-set construction.

\section{A one-input comparison at every depth}
In this section the fixed augmented input vector $u=(1,x)$ obeys $\|u\|_2^2=d$. The finite-set theorem uses the unit-sphere inputs in \eqref{eq:unit-input-set}, whose augmented squared norm is two; its separate construction starts in Section~\ref{sec:sphere-direct-kernel}.
\begin{theorem}[A finite-width one-output comparison]
\label{thm:main}
There is an absolute constant $C$ such that, for every $H\ge1$, $N\ge8$, and prescribed output $a\in\R$, there is a law $\rho_{H,N,a}\ll\mu$ with
\begin{align}
 \frac2N\KL(\rho_{H,N,a}\|\mu)&\le C(1+a^2)+C\frac HN,\label{eq:entropy}\\
 \E_\rho z&=a,\qquad \E_\rho(z-a)^2\le\frac1N.\label{eq:error}
\end{align}
The input dimension and chosen input do not enter $C$. In particular the normalized entropy cost is bounded independently of depth whenever $N\ge H$.
\end{theorem}

The comparison law need not have independent weights. Its conditionals are explicitly Gaussian, and its relative entropy is computed against the unchanged independent Gaussian prior. This is a valid comparison distribution for a Bayesian variational formula; it is not a modification of the learner's prior.

\section{Two scalar estimates}
Put
\[
 F(q)=\E\tanh^2(\sqrt q\,Z),\qquad Z\sim\mathcal N(0,1).
\]

\begin{lemma}[Critical variance drift]
\label{lem:drift}
For an absolute $c>0$ and every $0<q\le1/64$,
\[
 q-2q^2\le F(q)\le q-cq^2.
\]
The function $F$ is continuous and strictly increasing on $(0,\infty)$.
\end{lemma}
\begin{proof}
Monotonicity and continuity follow from the same properties of $\tanh^2(\sqrt q\,z)$ as a function of $q$ for every $z\ne0$, and bounded convergence. The elementary inequality
$\tanh^2 t\ge t^2-(2/3)t^4$ gives the lower bound after Gaussian averaging. To verify the inequality, when its right side is nonpositive it is immediate. Otherwise $|t|\le\sqrt{3/2}$ and $\tanh t\ge t-t^3/3$ for $t\ge0$, obtained by integrating $\tanh'(t)\ge1-t^2$; squaring gives the claim.

For $0\le t\le1$, concavity gives $\tanh t\ge t/2$. Hence
$t-\tanh t=\int_0^t\tanh^2 s\,ds\ge t^3/12$, and
$t^2-\tanh^2t\ge t^4/12$.
Restricting the Gaussian expectation to $|Z|\le1$ gives
$q-F(q)\ge q^2\E[Z^4\mathbf1_{|Z|\le1}]/12$.
This positive absolute constant is a valid $c$.
\end{proof}

\begin{lemma}[A uniform reciprocal empirical-variance bound]
\label{lem:inverse}
Let $0<v\le1/64$, let $Z_1,\ldots,Z_N$ be independent standard Gaussians, and put
\[
 Y=\frac1{NF(v)}\sum_{i=1}^N\tanh^2(\sqrt v Z_i).
\]
For every $N\ge8$, an absolute constant $C_1$ satisfies
\[
 \E Y=1,\qquad \E Y^{-2}\le C_1,\qquad
 \E(Y-1)^4\le\frac{C_1}{N^2},\qquad
 1\le\E Y^{-1}\le1+\frac{C_1}{N}.
\]
Moreover $\E\log Y\le0$, and this logarithm is integrable.
\end{lemma}
\begin{proof}
For $U=\tanh^2(\sqrt v Z)/F(v)$, Lemma~\ref{lem:drift} gives $U\le2Z^2$. Concavity of tanh on $[0,1]$ and monotonicity outside that interval give
\[
 U\ge\tfrac14\min\{Z^2,1\}.
\]
Thus all fixed positive moments of $U$ are uniformly bounded, and the fourth-moment formula for a mean of independent centered variables gives the stated bound for $Y-1$.

Write $W=\min\{Z^2,1\}$ and $f(t)=\E e^{-tW}$. There exists an absolute $b>0$ such that
\[
 f(t)\le(1+bt)^{-1/2}\qquad(t\ge0).
\]
Indeed, for $t>0$ the continuous positive function $(f(t)^{-2}-1)/t$ has positive limits $2\E W$ at zero and $2$ at infinity. The latter follows by integrating the Gaussian density near zero: $f(t)\sim(2t)^{-1/2}$. Its infimum is therefore positive. For independent copies $W_i$, the integral identity $x^{-2}=\int_0^\infty t e^{-tx}\,dt$ now yields
\[
 \E\Bigl(\sum_iW_i\Bigr)^{-2}
 \le\int_0^\infty t(1+bt)^{-N/2}\,dt
 =\frac{b^{-2}}{(N/2-1)(N/2-2)}.
\]
Since $Y\ge(4N)^{-1}\sum_iW_i$, this proves the reciprocal second-moment bound uniformly for $N\ge8$.

Finally the exact identity
$Y^{-1}=1-(Y-1)+(Y-1)^2/Y$ gives
\[
 \E Y^{-1}-1
 \le\bigl[\E(Y-1)^4\bigr]^{1/2}\bigl[\E Y^{-2}\bigr]^{1/2}
 \le C_1/N,
\]
after enlarging the constant. Jensen gives both $\E Y^{-1}\ge1$ and $\E\log Y\le0$. The positive and reciprocal moment bounds imply integrability of the logarithm.
\end{proof}

\section{The variance schedule and Gaussian control}
Fix $v_0=1/64$, set $q_{-1}=v_0$, and recursively set
$q_j=F(q_{j-1})$ for $j\ge0$. These positive numbers decrease to zero. Lemma~\ref{lem:drift} gives
\begin{equation}
 q_{j+1}\ge(31/32)q_j,\qquad
 cq_j^2\le q_j-q_{j+1}\le2q_j^2.
 \label{eq:schedule-bounds}
\end{equation}

For odd hidden depth $H=2M+1$, prescribe the sequence of mean hidden squared norms
\[
 q_0,q_1,\ldots,q_M,q_{M-1},\ldots,q_0.
\]
For even $H=2M+2$, append one further $q_0$ to this sequence. At any prescribed transition from mean norm $r$ to mean norm $r'$, let $v\le v_0$ be the unique value with $F(v)=r'$ and put $\alpha_0=v/r$.
On a forward transition, $\alpha_0=1$. On a reverse transition from $q_{j+1}$ to $q_j$,
\[
 \alpha_0=q_{j-1}/q_{j+1},\qquad 1\le\alpha_0\le2,
 \qquad \alpha_0-1\le5q_j.
\]
The possible repeated final transition also has $1\le\alpha_0=v_0/q_0\le2$.
Writing $h(\alpha)=\alpha-1-\log\alpha$, we have
$h(\alpha)\le(\alpha-1)^2/2$ for $\alpha\ge1$. Therefore the total deterministic transition cost satisfies
\begin{equation}
 \sum_{\text{transitions}} h(\alpha_0)
 \le C_2\sum_{j=0}^{M-1}q_j^2+C_2
 \le C_2\left(1+\frac{q_0}{c}\right),
 \label{eq:finite-action}
\end{equation}
independently of $H$. The last inequality telescopes~\eqref{eq:schedule-bounds}.

We now realize this schedule at finite width. In the first layer, change each row covariance from $I/d$ to
\[
 \frac1d\left[I+(v_0-1)\frac{uu^{\mathsf T}}{\|u\|^2}\right].
\]
The input preactivations are then independent $\mathcal N(0,v_0)$ variables. The first hidden norm has mean $q_0$.

At a later layer, let $h$ be the actual preceding hidden vector and let
$\widehat q=\|h\|^2/N$. Almost surely $\widehat q>0$. For a transition whose prescribed preactivation variance is $v$, change every row covariance from $I/N$ to
\begin{equation}
 \frac1N\left[I+(\alpha-1)\frac{hh^{\mathsf T}}{\|h\|^2}\right],
 \qquad \alpha=v/\widehat q.
 \label{eq:control}
\end{equation}
Given the past, the next preactivations are exactly independent $\mathcal N(0,v)$ variables. Their distribution does not depend on the past. In particular the next empirical norm divided by its prescribed mean has exactly the law in Lemma~\ref{lem:inverse}. All rows are centered, and all conditional covariance eigenvalues are positive.

For one Gaussian row, the relative entropy of this rank-one covariance change is $h(\alpha)/2$. Since $\widehat q=rY$, we have $\alpha=\alpha_0/Y$. Lemma~\ref{lem:inverse} gives
\begin{align}
 \E h(\alpha_0/Y)
 &=h(\alpha_0)+\alpha_0(\E Y^{-1}-1)+\E\log Y\notag\\
 &\le h(\alpha_0)+2C_1/N.
 \label{eq:feedback-cost}
\end{align}
The chain rule for relative entropy sums these conditional costs against the original independent Gaussian prior. The first layer contributes $Nh(v_0)/2$. Equations~\eqref{eq:finite-action}--\eqref{eq:feedback-cost} therefore bound the normalized hidden-weight entropy by an absolute constant plus $C H/N$.

\section{Readout and prior probability consequences}
At the final layer the empirical norm has mean $q_0$, regardless of $H$. Conditional on all hidden weights, give the readout its native covariance $I/N^2$ and mean
\[
 m_v=\frac{a h^H}{N\widehat q_H},\qquad
 \widehat q_H=\|h^H\|^2/N.
\]
Then $m_v^{\mathsf T}h^H=a$ exactly, while its conditional output variance is $\widehat q_H/N\le1/N$. Its conditional relative entropy is
\[
 \frac{N^2}2\|m_v\|^2=\frac{Na^2}{2\widehat q_H}.
\]
Lemma~\ref{lem:inverse} gives
$\E\widehat q_H^{-1}\le q_0^{-1}(1+C_1/N)$, an absolute constant. Combining this with the hidden-weight cost proves Theorem~\ref{thm:main}.

\begin{corollary}[One-output prior probabilities have no extra depth exponent]
\label{cor:probability}
There is an absolute $C$ such that, if $N\ge\max\{8,H,4/\varepsilon^2\}$, then for every $a\in\R$,
\[
 \mu\{|z-a|\le\varepsilon\}\ge\exp[-CN(1+a^2)].
\]
For $|a|>\varepsilon$, the corresponding elementary upper bound is
\[
 \mu\{|z-a|\le\varepsilon\}
 \le2\exp[-N(|a|-\varepsilon)^2/2].
\]
Thus at fixed nonzero output and sufficiently small fixed tolerance, negative log neighborhood probability is of order width, uniformly in depth whenever $N\ge H$.
\end{corollary}
\begin{proof}
Under the comparison law, Chebyshev gives probability at least $3/4$ to the neighborhood. The binary relative-entropy inequality implies
$\log(1/\mu(A))\le2\KL(\rho\|\mu)+2\log2$ whenever $\rho(A)\ge1/2$. Apply Theorem~\ref{thm:main} and absorb $H\le N$ into the constant.

Under the prior, conditional on hidden weights, the readout is centered Gaussian with output variance $\|h^H\|^2/N^2\le1/N$. Its Gaussian tail gives the second inequality.
\end{proof}

\section{A sharper small-amplitude comparison}
\label{sec:sharp-amplitude}
The constant initial variance shrink in Theorem~\ref{thm:main} is unnecessary when depth is long enough relative to the prescribed amplitude. Starting with the native first layer gives a comparison cost linear in small amplitude.

\begin{lemma}[The reciprocal bound on a larger variance interval]
\label{lem:inverse-large}
Lemma~\ref{lem:inverse} remains valid uniformly for $0<v\le1$, after changing its absolute constant.
\end{lemma}
\begin{proof}
The lower bound $\tanh^2(\sqrt v Z)\ge(v/4)\min\{Z^2,1\}$ holds for every $v\le1$. Hence $F(v)\ge c_0v$ for an absolute $c_0>0$, while $F(v)\le v$. The normalized variable $U$ is therefore at most $Z^2/c_0$ and at least $\min\{Z^2,1\}/4$. Every step of the earlier proof is unchanged.
\end{proof}

\begin{theorem}[Linear cost for a small prescribed output]
\label{thm:small-amplitude}
There are absolute constants $a_0,C_0>0$ such that, whenever
\[
 0<|a|\le a_0,\qquad H\ge C_0/|a|,\qquad N\ge8,
\]
there is a law $\rho\ll\mu$ satisfying
\begin{align}
 \frac2N\KL(\rho\|\mu)&\le C_0|a|+C_0H/N,\label{eq:linear-amplitude-cost}\\
 \E_\rho z&=a,\qquad \E_\rho(z-a)^2\le |a|/N.\label{eq:linear-amplitude-error}
\end{align}
The first layer retains its exact native Gaussian law.
\end{theorem}
\begin{proof}
For this proof put $q_{-1}=1$ and $q_j=F(q_{j-1})$. The upper drift estimate $F(q)\le q-cq^2$ is valid for all $q\le1$, by the same restriction to $|Z|\le1$ used in Lemma~\ref{lem:drift}. Thus
\[
 \frac1{q_{j+1}}\ge\frac1{q_j}+c,
 \qquad q_j\le\frac1{1/q_0+cj}.
\]
The sequence enters $(0,1/64]$ after an absolute number of steps. Choose $a_0>0$ sufficiently small that the first index $m$ with $q_m\le|a|$ occurs after this entry. The lower drift estimate then gives
\[
 |a|/2\le q_m\le|a|,\qquad m\le C/|a|.
\]
Increase $C_0$ so the depth hypothesis ensures $H\ge m+1$.

Let $M=\lfloor(H+m-1)/2\rfloor$. Use the mean-variance schedule
\[
 q_0,q_1,\ldots,q_M,q_{M-1},\ldots,q_m,
\]
appending one repeated $q_m$ if needed to obtain exactly $H$ hidden layers. The first layer is native, with preactivation variance one. Every forward transition has deterministic covariance factor one. Every reverse transition lies in the small-variance regime and has the factor analyzed in~\eqref{eq:finite-action}. Consequently its total deterministic cost is at most
\[
 C\sum_{j=m}^{M-1}q_j^2+Cq_m^2\le Cq_m\le C|a|.
\]
The possible repeated endpoint has cost $O(q_m^2)$, because its factor is $q_{m-1}/q_m=1+O(q_m)$.

Use the same exact feedback covariances~\eqref{eq:control}. Lemma~\ref{lem:inverse-large} makes the total finite-width correction at most $CH/N$, including the early forward transitions. The final conditional readout shift has normalized entropy
\[
 a^2\E\widehat q_H^{-1}\le\frac{a^2}{q_m}(1+C/N)\le C|a|.
\]
Its conditional output mean is exactly $a$, and its average output variance is $q_m/N\le|a|/N$. This proves both assertions.
\end{proof}

\section{Matching upper tails for the unchanged prior}
\label{sec:matching-upper}
Under the original independent Gaussian prior, the empirical hidden squared norms
$q_\ell=\|h^\ell\|^2/N$ form a Markov chain. Conditional on $q_\ell=q$,
\[
 q_{\ell+1}\ \stackrel{d}=\ \frac1N\sum_{i=1}^N\tanh^2(\sqrt q\,Z_i).
\]
They lie in $[0,1]$, and the first one is also at most one.

\begin{lemma}[An exponential supermartingale with quantitative drift]
\label{lem:variance-supermartingale}
There are absolute constants $\theta,b>0$ such that
\begin{equation}
 \E[e^{\theta N q_{\ell+1}}\mid q_\ell=q]
 \le e^{\theta Nq-bNq^2}\qquad(0\le q\le1).
 \label{eq:variance-supermartingale}
\end{equation}
In particular $e^{\theta Nq_\ell}$ is a positive supermartingale.
\end{lemma}
\begin{proof}
Put $U=\tanh^2(\sqrt qZ)$. We have $U\le qZ^2$, and for every integer $k\ge2$,
\[
 \E|U-\E U|^k\le2^k\E U^k
 \le4^kq^k k!.
\]
Expanding the exponential series gives, whenever $|t|q\le1/8$,
\[
 \log\E e^{t(U-\E U)}\le32t^2q^2.
\]
The global upper drift estimate $F(q)\le q-cq^2$ holds on $[0,1]$. Choose
$0<\theta\le\min\{1/8,c/64\}$ and set $b=\theta c/2$. Independence of the $N$ rows yields
\[
 \E e^{\theta Nq_{\ell+1}}
 \le\exp\{N[\theta F(q)+32\theta^2q^2]\}
 \le\exp\{\theta Nq-bNq^2\}.
\]
The case $q=0$ is immediate.
\end{proof}

\begin{lemma}[Decay to a small variance with exponential confidence]
\label{lem:variance-tail}
There are absolute constants $C,c>0$ such that for $0<r\le1/4$,
\begin{equation}
 H\ge C/r\quad\Longrightarrow\quad
 \Pr_\mu(q_H\ge r)\le C e^{-cNr}.
 \label{eq:variance-tail}
\end{equation}
\end{lemma}
\begin{proof}
First consider a chain started with $q\le R$, where $0<R\le1$. If it already lies below $R/2$, nothing is required. Otherwise stop at the first time it reaches $q\le R/2$ or $q\ge2R$. Before stopping, $q\ge R/2$, so~\eqref{eq:variance-supermartingale} shows that
\[
 \exp\{\theta N q_{t\wedge\tau}+(bNR^2/4)(t\wedge\tau)\}
\]
is a supermartingale. On $\{\tau>T\}$, it is at least
$\exp\{\theta NR/2+bNR^2T/4\}$. Its initial value is at most $e^{\theta NR}$, so
\[
 \Pr(\tau>T)\le\exp\{\theta NR/2-bNR^2T/4\}.
\]
The positive supermartingale $e^{\theta Nq_t}$ also gives probability at most $e^{-\theta NR}$ of reaching $2R$ before the lower endpoint. This follows by optional stopping at a finite horizon and then letting the horizon increase. Consequently, for an absolute $K$,
\begin{equation}
 \Pr\{\text{the chain has not reached }R/2
              \text{ within }\lceil K/R\rceil\text{ steps}\}
 \le2e^{-c_1NR}.
 \label{eq:phase-descent}
\end{equation}
The same bound holds conditional on any starting stopping time with $q\le R$.

Apply~\eqref{eq:phase-descent} at successive dyadic levels
$R_j=2^{-j}$, starting from $q_1\le1$. Stop when $R_J\le r/4$ for the first time. Then $R_J>r/8$, and the total allowed time is at most
$\sum_{j<J}\lceil K/R_j\rceil\le C/r$.
The probability any phase fails is bounded by
\[
 2\sum_{j<J}e^{-c_1NR_j}\le C_2e^{-c_2Nr}
\]
when $Nr\ge1$: the levels are geometrically spaced and the last starting level exceeds $r/4$. If $Nr<1$, the desired assertion follows by increasing $C_2$.

After reaching $q\le r/4$, the supermartingale in Lemma~\ref{lem:variance-supermartingale} bounds the probability of ever returning to $q\ge r$ by $e^{-3\theta Nr/4}$. Combining this with the phase bound proves~\eqref{eq:variance-tail} at every later time, after adjusting the absolute constant in the depth requirement to include the first layer.
\end{proof}

\begin{theorem}[Laplace-order one-input tails at critical scaling]
\label{thm:laplace-tails}
There are absolute $a_0,c,C>0$ such that, for $0<a\le a_0$ and $H\ge C/a$,
\begin{equation}
 \Pr_\mu(|z|\ge a)\le C e^{-cNa}.
 \label{eq:laplace-upper}
\end{equation}
If in addition $N\ge C H/a$, then for either sign of $a$,
\begin{equation}
 e^{-CN|a|}\le
 \mu\{|z-a|\le|a|/4\}\le C e^{-cN|a|}.
 \label{eq:laplace-neighborhood}
\end{equation}
Thus the precision-normalized one-output neighborhood cost is of order amplitude, rather than amplitude squared times depth, in this critical deep regime.
\end{theorem}
\begin{proof}
On $q_H\le a$, the conditional Gaussian readout has variance at most $a/N$, so
\[
 \Pr(|z|\ge a)\le\Pr(q_H>a)+2e^{-Na/2}.
\]
Lemma~\ref{lem:variance-tail} proves~\eqref{eq:laplace-upper}.

For the lower neighborhood bound use Theorem~\ref{thm:small-amplitude}. Its output variance is at most $|a|/N$, so the comparison law puts probability at least $3/4$ on the radius-$|a|/4$ neighborhood once $N|a|\ge64$. Its normalized entropy is at most $C|a|+CH/N$. The width hypothesis absorbs the second term into the first, and the binary relative-entropy inequality gives the lower bound. The upper bound follows from~\eqref{eq:laplace-upper} with threshold $3|a|/4$, after enlarging the depth constant. Negative prescribed outputs are covered by symmetry.
\end{proof}

\begin{corollary}[Small second moment and larger noncentral tails]
The unchanged prior satisfies $\E z^2\le C/(NH)$ for an absolute constant $C$. This is compatible with Theorem~\ref{thm:laplace-tails}: its amplitude regime has $H|a|$ bounded below, and its tail exponent is of order $N|a|$, rather than the order $NHa^2$ suggested by a Gaussian with that variance.
\end{corollary}
\begin{proof}
Put $m_\ell=\E q_\ell$. The global drift and Jensen give
$m_{\ell+1}\le m_\ell-cm_\ell^2$, so $m_H\le C/H$ by taking reciprocals. Conditional on hidden weights, the centered readout has variance $q_H/N$. Averaging proves the claim. The displayed comparison of exponents concerns the noncentral tail regime; it is not a rejection of a central-limit approximation on its own scale.
\end{proof}

\section{Sharp exponential moments and a finite-set entropy floor}
\label{sec:sharp-mgf}
A heat-equation calculation sharpens the tail constant and gives an entropy inequality simultaneously on any finite set of inputs. This upper bound does not require the inputs to have equal norm.

\begin{lemma}[A scalar heat supersolution]
\label{lem:heat}
For $3/2\le\theta<2$, there is an absolute $c_*>0$ such that, for every $0\le q\le1$,
\begin{equation}
 \log\E\exp\{\theta\tanh^2(\sqrt qZ)\}
 \le\theta q-\theta\kappa_\theta q^2,
 \qquad \kappa_\theta=(2-\theta)c_*e^{-2}.
 \label{eq:heat-bound}
\end{equation}
The constant can be chosen so that $0<\kappa_\theta<1/2$.
\end{lemma}
\begin{proof}
Put $f_\theta(x)=e^{\theta\tanh^2x}$ and $u=\tanh^2x$. Direct differentiation gives
\[
 \frac{f_\theta''(x)}{2f_\theta(x)}
 =\theta\{1+(2\theta-4)u+(3-4\theta)u^2+2\theta u^3\}.
\]
The quadratic polynomial
$A(u)=(2\theta-4)+(3-4\theta)u+2\theta u^2$ is convex on $[0,1]$. Its endpoint values are $-2(2-\theta)$ and $-1$. Since $\theta\ge3/2$,
\[
 A(u)\le-2(2-\theta),\qquad
 f_\theta''/2\le\theta f_\theta[1-2(2-\theta)u].
\]
Let $M_\theta(q)=\E f_\theta(\sqrt qZ)$. The heat identity gives
$M_\theta'(q)=\tfrac12\E f_\theta''(\sqrt qZ)$. All derivatives involved are bounded, so differentiation is justified. First $M_\theta'\le\theta M_\theta$ yields $M_\theta(q)\le e^{\theta q}$. Also $F(q)\ge c_*q$ on $[0,1]$, by the lower bound using $\min\{Z^2,1\}$ in Lemma~\ref{lem:inverse-large}. Since $f_\theta\ge1$,
\[
 \frac{d}{dq}\log M_\theta(q)
 \le\theta-2\theta(2-\theta)\frac{F(q)}{M_\theta(q)}
 \le\theta-2\theta(2-\theta)c_*e^{-2}q.
\]
Integrate from zero to $q$. Decreasing $c_*$ if necessary ensures the stated bound on $\kappa_\theta$.
\end{proof}

\begin{theorem}[Exponential moments at every depth]
\label{thm:sharp-mgf}
For the native variance chain and $3/2\le\theta<2$,
\begin{equation}
 \log\E_\mu e^{\theta Nq_H}
 \le\frac{\theta N}{1+\kappa_\theta(H-1)}.
 \label{eq:sharp-variance-mgf}
\end{equation}
For $0<t<2$, put
\[
 \theta_t=\max\{3/2,t^2/2\},\qquad
 D_H(t)=\frac{t^2/2}{1+\kappa_{\theta_t}(H-1)}.
\]
Then for every input $x$,
\begin{equation}
 \log\E_\mu e^{tNz(x)}\le N D_H(t),\qquad
 \log\E_\mu e^{-tNz(x)}\le N D_H(t).
 \label{eq:sharp-output-mgf}
\end{equation}
For each fixed $t<2$, $D_H(t)=O_t(1/H)$.
\end{theorem}
\begin{proof}
Lemma~\ref{lem:heat} and conditional row independence give
\[
 \E[e^{\theta Nq_{\ell+1}}\mid q_\ell=q]
 \le e^{\theta Nq-\theta\kappa_\theta Nq^2}.
\]
The function
\[
 g(y)=y\exp[-\beta(\log y)^2],\qquad
 \beta=\frac{\kappa_\theta}{\theta N},
\]
is concave for $1\le y\le e^{\theta N}$. Indeed its second derivative has the sign of
$2\beta(\log y)^2-\log y-1$, which is negative because
$2\beta\log y\le2\kappa_\theta<1$.
Writing $M_\ell=\E e^{\theta Nq_\ell}$, Jensen therefore gives
$M_{\ell+1}\le g(M_\ell)$. Thus
\[
 m_\ell:=\frac{\log M_\ell}{\theta N}
 \quad\text{satisfies}\quad
 m_{\ell+1}\le m_\ell-\kappa_\theta m_\ell^2,
 \qquad m_1\le1.
\]
Taking reciprocals proves~\eqref{eq:sharp-variance-mgf}. Conditional on hidden weights, $z(x)$ is a centered Gaussian with variance $q_H(x)/N$, so its exponential moment equals the variance-chain moment with parameter $t^2/2$. If $t^2/2<3/2$, apply Jensen or H\"older to the larger parameter $\theta_t$. This proves~\eqref{eq:sharp-output-mgf}. The first hidden norm is always at most one, so the bound is uniform over inputs without an input-norm assumption.
\end{proof}

\begin{corollary}[A sharp finite-set entropy bound]
\label{cor:sup-entropy}
For a finite set $\mathcal A$ of $m\ge1$ inputs, let
$\|z\|_{\infty,\mathcal A}=\max_{x\in\mathcal A}|z(x)|$. Every $\rho\ll\mu$ and $0<t<2$ satisfy
\begin{equation}
 \frac2N\KL(\rho\|\mu)
 \ge2t\E_\rho\|z\|_{\infty,\mathcal A}
       -2D_H(t)-\frac{2\log(2m)}N.
 \label{eq:sup-entropy}
\end{equation}
The expectation may be interpreted by truncation. Thus, as depth and width grow with $\log m=o(N)$, the entropy penalty has coefficient arbitrarily close to four in front of the expected maximum output magnitude.
\end{corollary}
\begin{proof}
The exponential of a maximum is at most the sum of exponentials. Using both signs in~\eqref{eq:sharp-output-mgf},
\[
 \E_\mu e^{tN\|z\|_{\infty,\mathcal A}}
 \le2m e^{ND_H(t)}.
\]
The entropy variational inequality gives~\eqref{eq:sup-entropy}, first for truncations and then by monotone convergence.
\end{proof}

For a Gaussian-likelihood posterior with free energy
$F_S=\E_\Pi\widehat R_S+(2\lambda/N)\KL(\Pi\|\mu)$,
Corollary~\ref{cor:sup-entropy} gives the exact objective inequality
\[
 \E_\Pi[\widehat R_S+2t\lambda\|z\|_{\infty,S}]
 \le F_S+2\lambda D_H(t)+\frac{2\lambda\log(2n)}N.
\]
This is an entropy lower bound. Turning it into a sample-complexity conclusion would also require a matching function comparison on the full input domain.

\section{The leading constant in the one-output comparison}
\label{sec:sharp-comparison}
The partial variance recovery also gives the matching leading constant.

\begin{theorem}[Asymptotically sharp amplitude cost]
\label{thm:sharp-comparison}
There are absolute constants $a_0,C>0$ such that for
$0<|a|\le a_0$, $H\ge C/|a|$, and $N\ge8$, the comparison law can be chosen with
\begin{equation}
 \frac2N\KL(\rho\|\mu)
 \le4|a|+Ca^2+CH/N,
 \qquad \E_\rho z=a,\qquad
 \E_\rho(z-a)^2\le\frac{|a|}{2N}.
 \label{eq:sharp-comparison}
\end{equation}
Conversely, for every $\varepsilon>0$ there is $C_\varepsilon<\infty$ such that every law with $\E_\rho z=a$ satisfies
\begin{equation}
 \frac2N\KL(\rho\|\mu)
 \ge(4-\varepsilon)|a|-C_\varepsilon/H-2\log2/N.
 \label{eq:sharp-comparison-lower}
\end{equation}
\end{theorem}
\begin{proof}
Gaussian averaging of the Taylor expansion of $\tanh^2$ gives
\[
 F(q)=q-2q^2+O(q^3)\qquad(q\downarrow0),
\]
with an absolute uniform remainder on a small fixed interval. One elementary justification is
$\tanh^2x=x^2-(2/3)x^4+O(x^6)$ globally, with a larger constant outside $[-1,1]$, followed by the Gaussian sixth moment.

Use the native-first-layer schedule of Theorem~\ref{thm:small-amplitude}, now stopping the reverse trajectory at the first $q_m\le|a|/2$. Then
\[
 q_m=\frac{|a|}{2}(1+O(|a|)),\qquad q_m\le|a|/2.
\]
For a reverse transition indexed by $j$,
\[
 \alpha_j=\frac{q_{j-1}}{q_{j+1}}=1+4q_j+O(q_j^2),\qquad
 h(\alpha_j)=8q_j^2+O(q_j^3)
            =4(q_j-q_{j+1})+O(q_j^3).
\]
Since $\sum_{j\ge m}q_j^3\le Cq_m^2$, its total deterministic reverse cost is at most $4q_m+Cq_m^2$. A possible repeated endpoint costs $O(q_m^2)$. The feedback contribution remains $CH/N$.

The readout costs
\[
 \frac{a^2}{q_m}(1+C/N)
 \le2|a|+Ca^2+C|a|/N.
\]
Combining it with $4q_m\le2|a|$ proves the entropy upper bound, and the readout variance is $q_m/N\le|a|/(2N)$. The term $C|a|/N$ is absorbed into $CH/N$.

For the lower bound apply~\eqref{eq:sup-entropy} to one input, choose $t<2$ with $2t\ge4-\varepsilon$, and use $\E|z|\ge|\E z|=|a|$. The bound $D_H(t)=O_t(1/H)$ proves~\eqref{eq:sharp-comparison-lower}.
\end{proof}

In particular, the leading precision-normalized entropy cost tends to $4|a|$ when $a\to0$, $H|a|\to\infty$, and $H/(N|a|)\to0$. The corresponding prior large-deviation exponent is $2N|a|$. Theorem~\ref{thm:laplace-tails} states a two-sided neighborhood version with absolute constants; the sharp moment and comparison estimates above give leading constants arbitrarily close to two after the neighborhood radius is made a sufficiently small fraction of $|a|$ and width is increased accordingly.

\section{A quantitative feature kernel on the unit sphere}
\label{sec:sphere-direct-kernel}
Use the input set $\mathcal U_D$ from~\eqref{eq:unit-input-set}.
Let $A$ have independent $N(0,1/d)$ coordinates in $\R^{d+1}$. For $0\le\nu\le1$, define
\[
 \phi_A(x)=\E_G\tanh\left(A\cdot(1,x)+\sqrt{2\nu/d}\,G\right),
 \qquad K_\nu(x,x')=\E_A\phi_A(x)\phi_A(x').
\]
All finite-set inner products use the uniform probability law on $\mathcal U_D$.

\begin{lemma}[A small-variance Gaussian feature kernel]
\label{lem:sphere-direct-kernel}
The kernel operator $K_\nu$ has smallest eigenvalue at least
\[
 \widetilde\Lambda_D=e^{-40(D+1)}
\]
for every $D\ge1$ and every $0\le\nu\le1$.
\end{lemma}
\begin{proof}
Write $\alpha=\sqrt{2/d}$. The first Gaussian mean preactivation has variance $\alpha^2$ at every input. Expand
\[
 t\longmapsto\E_G\tanh(\alpha t+\alpha\sqrt\nu G)
\]
in orthonormal Hermite polynomials for a standard Gaussian variable. Even coefficients vanish. For odd $j\ge1$, the identities
\[
 \tanh t=\int_0^\infty\frac{\sin(\omega t)}{\sinh(\pi\omega/2)}\,d\omega,
 \qquad \E[\operatorname{He}_j(Z)e^{i\omega Z}]=(i\omega)^j e^{-\omega^2/2}
\]
give the exact absolute coefficient below. The Fourier identity follows by expanding $1/\sinh(\pi\omega/2)=2\sum_{k\ge0}e^{-(k+1/2)\pi\omega}$, integrating each sine term, and using the partial-fraction expansion of tanh. Gaussian decay justifies the Hermite calculation, and the integrands are integrable at zero. The coefficient is
\[
 |b_j|=\frac1{\alpha\sqrt{j!}}\int_0^\infty
 \frac{\omega^j e^{-(1+\nu)\omega^2/2}}
      {\sinh(\pi\omega/(2\alpha))}\,d\omega.
\]
For $j\le D+4=d+2$, restrict to $\omega\in[\sqrt j/2,\sqrt j]$. Since $\nu\le1$, $\sqrt{j!}\le j^{j/2}$, and $1/\sinh u\ge2e^{-u}$,
\[
 |b_j|\ge\frac{\sqrt j}{\alpha}
       \exp\{-j\log2-j-(\pi/2)\sqrt{jd/2}\}\ge e^{-5d}.
\]
For the final inequality use $\sqrt j/\alpha\ge1$, $j\le d+2$, $\sqrt{jd/2}\le d$ for $d\ge2$, and
$(1+\log2)(d+2)+(\pi/2)d\le5d$ for $d\ge3$.

The correlation of two normalized first preactivations is
\[
 c(x,x')=\frac{1+TT'}2+\frac{z\cdot z'}{8D}
            +\frac{\sqrt{3/4-T^2}\sqrt{3/4-(T')^2}}2.
\]
Thus $K_\nu=\sum_{j\ \mathrm{odd}}b_j^2c^j$, with every multinomial term a positive semidefinite rank-one kernel. For a nuisance subset $S$ of size $s$ and $a\in\{0,1,2,3\}$, choose $b\in\{0,1\}$ so that $j=s+a+b$ is odd. Retaining the term $T^a\chi_S(z)\otimes T^a\chi_S(z)$ gives coefficient
\[
 b_j^2\frac{j!}{a!b!}\,2^{-(a+b)}(8D)^{-s}
 \ge \frac1{96}e^{-10d}\frac{s!}{(8D)^s}
 \ge\frac1{96}e^{-10d-8D}.
\]
The last inequality follows because $(8D)^s/s!$ is one term of the series for $e^{8D}$. All selected terms are distinct. The four-point Gram matrix of $1,T,T^2,T^3$ has determinant $9/65536$ and trace less than two: the unnormalized Vandermonde determinant is $3/16$. Its minimum eigenvalue is therefore at least $9/524288>2^{-16}$. The nuisance Walsh characters are orthonormal. Consequently
\[
 \lambda_{\min}(K_\nu)
 \ge\frac{2^{-16}}{96}e^{-18D-20}\ge e^{-40(D+1)}.
\]
This proves the claim.
\end{proof}


\section{A coherent seed at quadratic logarithmic-domain cost}
\label{sec:canonical-spike-prefix}
Write $F(q)=\E\tanh^2(\sqrt qG)$ and $\chi(q)=1-F(q)$. A covariance spike of size $1+\alpha$ evolves variance and directional correlation by
\begin{equation}
 v=q+\alpha r^2,\qquad q'=F(v),\qquad
 r'=\sqrt{1+\alpha}\,r\chi(v).
 \label{eq:canonical-spike-update}
\end{equation}
For a valid state $0<r^2\le q\le1$, its coherence is $\theta=r^2/q$.

\begin{lemma}[Native critical decay preserves coherence up to a constant]
\label{lem:native-coherence}
There is an absolute $C$ such that any number of zero-spike transitions satisfy $\theta_j\ge e^{-C}\theta_1$.
\end{lemma}
\begin{proof}
For a zero spike the coherence ratio is
\[
 L(q)=\frac{q\chi(q)^2}{F(q)}.
\]
Stein and Cauchy--Schwarz give $0<L(q)\le1$. Taylor expansion at zero gives
\[
 F(q)=q-2q^2+\tfrac{17}{3}q^3+O(q^4),\qquad
 L(q)=1-\tfrac23q^2+O(q^3).
\]
Thus $-\log L(q)/(q-F(q))$ has finite limit $1/3$ at zero. It is continuous and finite on $(0,1]$, so $-\log L(q)\le C(q-F(q))$ there. The variance follows $q'=F(q)$; summing this inequality telescopes to at most $Cq_1\le C$.
\end{proof}

\begin{lemma}[A constant-coherence prefix]
\label{lem:spike-prefix}
Fix $C_0,c_0>0$. There are absolute positive constants $K,C,c$ depending only on them such that the following holds for all sufficiently large $H$ and $1\le D\le\sqrt H$. Suppose
\[
 c_0/(D+1)\le q_1\le1,\qquad
 e^{-C_0(D+1)}\le|r_1|\le\sqrt{q_1}/2.
\]
Set $\alpha=K(D+1)/H$. Use zero spikes for the first $\lfloor H/4\rfloor$ transitions and then constant spike $\alpha$. At every step between $\lceil3H/4\rceil$ and $H$,
\[
 c\alpha\le q_\ell\le C\alpha,
 \qquad r_\ell^2/q_\ell\ge1/2.
\]
The covariance entropy action is at most $K^2(D+1)^2/(2H)$. The fixed constant $K$ may be enlarged as needed, after increasing the depth threshold.
\end{lemma}
\begin{proof}
Native reciprocal variance estimates give $c/H\le q\le C/H$ at the end of the initial burn, using $q_1\ge c_0/(D+1)$ and $D\le\sqrt H$. Lemma~\ref{lem:native-coherence} gives coherence at least $e^{-C_1(D+1)}$ there; no logarithmic loss in $H$ occurs.

For the subsequent constant spike, $r^2\le q$ and the quadratic drift imply
\[
 q'\le F((1+\alpha)q)\le(1+\alpha)q-cq^2.
\]
Starting from $q\le C/H$, this preserves $q\le C_2\alpha$ after choosing a large enough fixed $C_2$. The native lower comparison $q'\ge F(q)$ also ensures $q\ge c/H$ over the full horizon.

The exact coherence update is
\[
 \theta'=\frac{(1+\alpha)\theta}{1+\alpha\theta}
             L(q(1+\alpha\theta)).
\]
Since $L(v)\ge1-Cv^2$ at small $v$ and $q\le C_2\alpha$, while $\theta\le1/2$ this gives $\theta'\ge(1+\alpha/4)\theta$ for small enough $\alpha$. It reaches $1/2$ within $C_3(D+1)/\alpha=C_3H/K$ steps. Thereafter the boundary $1/2$ is invariant, because the logistic increment there is of order $\alpha$ and dominates the order-$\alpha^2$ loss.

Once $\theta\ge1/2$, whenever $q\le c_1\alpha$ for sufficiently small $c_1>0$,
\[
 q'=F(q(1+\alpha\theta))\ge q(1+\alpha/4).
\]
It reaches $c_1\alpha$ from $c/H$ within
$C_4\alpha^{-1}\log(C_5\alpha H)$ steps. This is at most
$C_6H\log(C_7K)/K$, uniformly in $D\ge1$, because
$\log(C_5K(D+1))/(D+1)\le C\log(C_7K)$.
The lower boundary is invariant by monotonicity of $F$. Choose the fixed $K$ so that the two growth periods total less than $H/4$. The claimed state bounds then hold by step $3H/4$ and remain valid. Finally $\alpha-\log(1+\alpha)\le\alpha^2/2$ proves the action bound.
\end{proof}

\begin{lemma}[Exact orthogonal-defect identity]
\label{lem:spike-defect}
For any nonnegative spike, put $e=q-r^2$. Then
\[
 e'=J(q')+e(1-q')^2,
 \qquad J(u)=u-F^{-1}(u)(1-u)^2.
\]
For small $u\ge0$, $J(u)=\tfrac23u^3+O(u^4)$ and $0\le J(u)\le Cu^3$.
\end{lemma}
\begin{proof}
Since $(1+\alpha)r^2=v-e$ and $q'=F(v)$, the displayed recurrence gives $r'^2=(v-e)(1-q')^2$. Taylor expansion yields $F^{-1}(u)=u+2u^2+(7/3)u^3+O(u^4)$ and hence the expansion for $J$. Nonnegativity follows from $v(\E\tanh'(\sqrt vG))^2\le\E\tanh^2(\sqrt vG)$.
\end{proof}

\section{A sharp reverse tail from constant coherence}
\label{sec:covariance-spike-tail}
The reverse phase does not require an almost perfectly coherent seed. A fixed positive fraction of coherent variance suffices.

\begin{lemma}[Reverse tail with a controlled initial defect]
\label{lem:spike-reverse-tail}
Fix $0<\rho_0<1$ and $B_0<\infty$. There are $q_*>0$ and $C<\infty$, depending only on these constants, with the following property. Suppose
\[
 0<q_0\le Q\le q_*,\qquad
 0\le e_0\le\rho_0q_0,\qquad r_0=\sqrt{q_0-e_0}>0.
\]
There is a deterministic sequence of covariance controls $1\le R^2\le2$, using at most $C/q_0$ reverse steps, which ends with
\[
 Q-CQ^2\le q_f\le Q,\qquad
 e_f\le e_0+CQ^2,\qquad e_f/q_f\le\rho_1<1,
\]
where $\rho_1$ is a fixed constant depending only on $\rho_0$, and whose action satisfies
\[
 \sum h(R^2)\le4(Q-q_0)+Ce_0+CQ^2,
 \qquad h(t)=t-1-\log t.
\]
Before these reverse steps, up to $B_0/q_0$ prescribed padding steps can hold the variance at $q_0$. They add action at most $Cq_0$ and change the defect upper bound to $e_0+Cq_0^2$.
\end{lemma}
\begin{proof}
Choose $\rho_1\in(\rho_0,1)$ and decrease $q_*$ as needed. For a reverse step prescribe $q'=F^{-1}(q)$ and use
\[
 R^2-1=\frac{F^{-1}(q')-q}{q-e}.
\]
Writing $\rho=e/q$, Taylor expansion gives uniformly for $\rho\le\rho_1$,
\[
 q'-q=2q^2+O(q^3),\qquad
 R^2-1=\frac{4q}{1-\rho}+O(q^2).
\]
Thus the spike is positive and at most one for small $q_*$. The defect identity implies
\[
 \rho'=\frac{J(q')}{q'}+\rho\frac q{q'}(1-q')^2
 \le C(q')^2+\rho(1-q').
\]
At $\rho=\rho_1$, this points inward for sufficiently small $q'$, so $\rho\le\rho_1$ remains invariant. Also $e'\le e+C(q')^3$.

The action expansion is
\[
 h(R^2)=\frac{8q^2}{(1-\rho)^2}+O(q^3)
 \le4(q'-q)+C\rho q^2+Cq^3.
\]
The extra term telescopes through the exact defect:
\[
 e-e'=e\{1-(1-q')^2\}-J(q')
       \ge e q'-C(q')^3.
\]
Consequently $\sum\rho q^2=\sum e q\le e_0+C\sum(q')^3$.
Run reverse steps until the next would exceed $Q$. As before, the last unused increment is $O(Q^2)$, the number of steps is at most $C/q_0$, and
\[
 \sum(q')^3\le C\sum q(q'-q)\le CQ^2.
\]
Summing the action and defect inequalities proves all claims for the reverse phase.

For padding at $q_0$, use
$R^2-1=[F^{-1}(q_0)-q_0]/(q_0-e)$. While $e/q_0\le\rho_1$, this is $O(q_0)$ and costs $O(q_0^2)$ per step. The defect recurrence keeps that ratio invariant at small $q_0$, and $e'\le e+Cq_0^3$. At most $B_0/q_0$ steps therefore add at most $Cq_0$ action and $Cq_0^2$ defect. The fixed gap between $\rho_0$ and $\rho_1$ allows the reverse argument to start after this padding, after a harmless enlargement of constants.
\end{proof}

\begin{corollary}[Leading cost four times the requested amplitude]
\label{cor:spike-four-amplitude}
Let $a>0$ be small and take $Q=a/2$ in Lemma~\ref{lem:spike-reverse-tail}. If $q_0\le Q$, choose readout coefficient $B=a/\sqrt{q_f-e_f}$. The endpoint output is exactly $a$, and covariance action plus readout normalized entropy is at most
\[
 \sum h(R^2)+B^2\le4a+Ca^2+Cq_0,
\]
including the allowed padding. Every target value in $[-a,a]$ can be obtained by prewarping a continuously and oddly parametrized input seed whose endpoint attains $a$.
\end{corollary}
\begin{proof}
Since $e_f/q_f\le\rho_1<1$,
\[
 B^2=\frac{a^2}{q_f-e_f}
 \le\frac{a^2}{q_f}+C\frac{a^2e_f}{q_f^2}
 \le2a+Ca^2+Ce_f.
\]
The defect bound gives $e_f\le e_0+Ca^2\le Cq_0+Ca^2$. The reverse action is at most $2a+Ca^2+Cq_0$, including padding. Summing proves the claim. Continuity, oddness, and the intermediate value theorem give the final assertion; monotonicity is unnecessary.
\end{proof}

% Research module: finite-width realization of covariance-spike scalar schedules.
\section{Finite-width realization and arbitrary finite functions}
\label{sec:spike-realization}
Throughout this section the prior is~\eqref{eq:canonical-prior}.

Let $\mathcal X$ be a finite input set of size $m$, all with the same augmented squared norm. Suppose the native first hidden-feature kernel, as an operator for the uniform input law, has minimum eigenvalue at least $\Lambda>0$. Write
\[
 b_1=\E_\mu\tanh^2(W^1_i(1,x)),\qquad
 \delta=\sqrt\Lambda/2.
\]
The value $b_1$ is independent of $x$. The unit-sphere sets $\mathcal U_D$ satisfy these assumptions with $m=4\cdot2^D$, $b_1=F(2/d)$, and $\Lambda=e^{-40(D+1)}$.

\begin{lemma}[A unit direction realizing an arbitrary first-layer seed]
\label{lem:unit-seed}
Fix $\varepsilon>0$. At every sufficiently large $N$, an event $E$ under the unchanged first-layer prior has probability at least $1/2$, on which
\[
 \sup_{x\in\mathcal X}|\|h^1(x)\|^2/N-b_1|\le\varepsilon,
\]
and, for every function $u:\mathcal X\to[-1,1]$, there is a measurable unit vector $e_1\in\R^N$ with
\[
 e_1^{\mathsf T}h^1(x)/\sqrt N=\delta u(x)
                     \quad(x\in\mathcal X).
\]
A sufficient width is
\[
 N\ge C(\Lambda^{-2}+\varepsilon^{-2})\log(8m)+2m.
\]
Conditioning the first layer on $E$ has normalized entropy cost at most $2\log2/N$.
\end{lemma}
\begin{proof}
Let $H_1$ be the $m$ by $N$ first-feature matrix. Hoeffding and a union bound give a sample kernel
$K_N=H_1H_1^{\mathsf T}/N$ whose uniform-measure integral operator is at least $\Lambda I/2$, and simultaneously give the displayed diagonal concentration, with joint probability at least $1/2$.
For a prescribed $u$, the minimum-norm solution of
$H_1e=\sqrt N\delta u$ has squared norm
\[
 \delta^2u^{\mathsf T}K_N^{-1}u
 =\delta^2\langle u,(K_N/m)^{-1}u\rangle_{\mathrm{unif}}
 \le\frac{2\delta^2}{\Lambda}\|u\|_{\mathrm{unif}}^2\le1/2.
\]
The equality uses the normalization of the integral operator. Since $N>m$, add a vector from the nullspace of $H_1$, orthogonal to the minimum-norm solution, to make the norm exactly one. A measurable choice follows by projecting the coordinate basis onto that nullspace and choosing the first nonzero projection, with deterministic tie breaking. Conditioning has entropy $\log(1/\Pr(E))\le\log2$.
\end{proof}

\begin{lemma}[An exact conditional Gaussian covariance construction]
\label{lem:spike-finite-width}
Fix a deterministic schedule $0\le\alpha_\ell\le1$ for $2\le\ell\le H$, a scalar readout factor $0\le B\le B_0$, and a seed function $u:\mathcal X\to[-1,1]$. Consider the scalar recursion at each input
\[
 q_1=b_1,\quad r_1=\delta u(x),\qquad
 v_\ell=q_{\ell-1}+\alpha_\ell r_{\ell-1}^2,
\]
\[
 q_\ell=F(v_\ell),\qquad
 r_\ell=\sqrt{1+\alpha_\ell}\,r_{\ell-1}(1-q_\ell).
\]
For any fixed tolerance $\tau>0$, all sufficiently large widths admit a comparison law $\rho$ such that
\begin{equation}
 \frac2N\KL(\rho\|\mu)
 \le\frac{2\log2}N+
      \sum_{\ell=2}^H\{\alpha_\ell-\log(1+\alpha_\ell)\}+B^2,
 \label{eq:spike-exact-entropy}
\end{equation}
and
\[
 \left(\E_\rho\max_{x\in\mathcal X}|z_w(x)-B r_H(x)|^2\right)^{1/2}\le\tau.
\]
When $\log m+\log(1/\Lambda)=O(H)$ and $B_0,\tau$ are fixed, a sufficient width is $N\ge e^{CH}$ for a constant depending only on those bounds and the constants in the preceding lemmas.
\end{lemma}
\begin{proof}
Condition the native first layer on the event in Lemma~\ref{lem:unit-seed}, choosing its accuracy below. Construct its unit direction $e_1$ for $u$.
At layer $\ell$, conditional on all preceding weights, draw its rows independently from the centered Gaussian law with covariance
\[
 \frac1N[I+\alpha_\ell e_{\ell-1}e_{\ell-1}^{\mathsf T}].
\]
Set
\[
 \xi_{\ell i}=\frac{\sqrt N}{\sqrt{1+\alpha_\ell}}
                   (W^\ell e_{\ell-1})_i,
 \qquad e_\ell=\xi_\ell/\|\xi_\ell\|.
\]
Conditionally the $\xi_{\ell i}$ are independent standard Gaussians. The normalization is well defined almost surely. Given a previous hidden vector, write its empirical variance and correlation as
$\widehat q=\|h\|^2/N$, $\widehat r=e_{\ell-1}^{\mathsf T}h/\sqrt N$. Each new preactivation has the joint representation
\[
 \sqrt{1+\alpha_\ell}\,\widehat r\,\xi_i
       +\sqrt{\widehat q-\widehat r^2}\,G_i,
\]
where $(\xi_i,G_i)$ are independent standard Gaussian pairs. Gaussian integration by parts gives its conditional output-square mean $F(\widehat q+\alpha_\ell\widehat r^2)$ and its $\xi_i$-weighted output mean
\[
 \sqrt{1+\alpha_\ell}\,\widehat r
       \{1-F(\widehat q+\alpha_\ell\widehat r^2)\}.
\]
These are exactly the stated scalar update.

At a fixed layer, conditional on the past, bounded-variable concentration for output squares, sub-Gaussian concentration for $\xi_i\tanh(\cdot)$, and chi-square concentration for $N^{-1}\sum_i\xi_i^2$ show that the empirical variance and correlation differ from these conditional update means by at most $C\varepsilon$, simultaneously at all $m$ inputs, with failure at most $Cm e^{-cN\varepsilon^2}$. The sub-Gaussian bound is uniform because $|\xi_i\tanh(\cdot)|\le|\xi_i|$; centering preserves an absolute sub-Gaussian parameter. The denominator $\|\xi\|/\sqrt N$ changes the correlation by at most another constant times $\varepsilon$ on the same event.

The scalar map is uniformly Lipschitz on $0\le r^2\le q\le1$ and $0\le\alpha\le1$. Indeed $|F'|\le1$ by Gaussian differentiation and the bounded second derivative of $\tanh^2$, and the displayed formulas have uniformly bounded partial derivatives there. One may take an absolute Lipschitz constant ten. A union bound and induction consequently make the empirical state within $C10^H\varepsilon$ of the scalar state at every input, with failure at most $CHm e^{-cN\varepsilon^2}$.
Choose $\varepsilon$ a fixed multiple of $\tau/[(1+B_0)10^H]$ and increase $N$ until this failure probability is at most $\tau^2/[8(1+B_0)^2]$.

Finally, conditionally on the hidden weights, draw the readout with its native covariance $I/N^2$ and mean $B e_H/\sqrt N$. Its conditional output mean is $B\widehat r_H$, and its variance is at most $1/N$ at every input. A Gaussian union bound and integration give conditional expected maximum squared readout noise at most $C\log(2m)/N$. Both empirical and scalar correlations have absolute value at most one. The common bounded-error event and its complement therefore give the claimed expected maximum squared error, after a harmless tightening of $\varepsilon$ and increase of $N$.

For entropy, each hidden row changes just one covariance eigenvalue from one to $1+\alpha_\ell$. Its relative entropy is
$\tfrac12\{\alpha_\ell-\log(1+\alpha_\ell)\}$, and there are $N$ independent rows. The conditional readout mean shift has entropy $NB^2/2$. Chain-rule relative entropy adds these quantities and the initial conditioning cost. All directions are measurable functions of previously drawn parameters, so no extra entropy term is introduced by their use. This proves~\eqref{eq:spike-exact-entropy}.
\end{proof}

\begin{corollary}[A scalar endpoint suffices for every bounded finite function]
\label{cor:spike-ivt}
Fix a deterministic spike schedule and let $r_H(u)$ denote its final scalar correlation when $q_1=b_1$ and $r_1=\delta u$. Suppose $r_H(1)>0$. For every $f:\mathcal X\to[-a,a]$, set $B=a/r_H(1)$. Then there is a seed function $u:\mathcal X\to[-1,1]$ for which $B r_H(u(x))=f(x)$ exactly. Thus Lemma~\ref{lem:spike-finite-width} gives a function memorizer with the scalar endpoint's covariance cost and readout cost $a^2/r_H(1)^2$.
\end{corollary}
\begin{proof}
The recursion is continuous in $u$, its variance is even in $u$, and its correlation is odd. Hence the continuous odd function $r_H(u)$ takes every value in $[-r_H(1),r_H(1)]$ by the intermediate value theorem. Choose one preimage at each of the finitely many input values. No monotonicity assertion is needed.
\end{proof}

% Research module: quantitative finite-width improvement for the canonical spike construction.
\section{Polynomial width in the finite support size}
\label{sec:canonical-polynomial-width}
The exponential-in-depth width in Lemma~\ref{lem:spike-finite-width} is convenient but unnecessary for the schedules used here. Their small spikes allow a sharper stability estimate.

\begin{theorem}[A quantitative width for the finite-function comparison]
\label{thm:canonical-polynomial-width}
In the finite-function comparison constructed from Lemma~\ref{lem:spike-prefix} and Lemma~\ref{lem:spike-reverse-tail}, let $1\le D\le\sqrt H$ and $m=4\cdot2^D$. There are absolute constants $C,c>0$, depending only on the fixed small amplitude cap and the constants in the construction, such that every prescribed root-mean-square supremum tolerance $0<\tau<1$ is achieved at every width
\[
 N\ge C\tau^{-C}H^C e^{CD}.
\]
The normalized entropy bound is unchanged. In particular the comparison width is polynomial in the support size, the depth, and the inverse accuracy.
\end{theorem}
\begin{proof}
We first bound propagation of scalar errors uniformly over every prewarped seed $u\in[-1,1]$. Write
\[
 \mathcal T_\alpha(q,r)=
 \bigl(F(q+\alpha r^2),\sqrt{1+\alpha}\,r[1-F(q+\alpha r^2)]\bigr).
\]
On a sufficiently small fixed variance interval, $0\le F'\le1$, and its Jacobian satisfies
\begin{align*}
 |\partial_q\mathcal T_\alpha^{(q)}|&\le1,
 &|\partial_r\mathcal T_\alpha^{(q)}|&\le2\alpha\sqrt q,\\
 |\partial_q\mathcal T_\alpha^{(r)}|&\le C\sqrt q,
 &|\partial_r\mathcal T_\alpha^{(r)}|&\le1+C\alpha.
\end{align*}
The final estimate follows directly from
$\partial_r\mathcal T_\alpha^{(r)}=\sqrt{1+\alpha}(1-F(v)-2\alpha r^2F'(v))$; for small $q$ this derivative is nonnegative and at most $\sqrt{1+\alpha}$. Enlarging the fixed constants accommodates a fixed larger initial variance interval.

During the initial native burn, $\alpha=0$. The variance update is independent of $r$, while both diagonal derivatives are at most one in absolute value. Products of these triangular Jacobians have norm at most $C(1+H)$: the lower off-diagonal entry is a sum of at most $H$ bounded terms. Thus this stage amplifies errors by only a polynomial in $H$.

For the constant-spike prefix put $a=K(D+1)/H$. Every seed, not just the endpoint seed, obeys $q\le C_1a$ after the burn. Indeed $r^2\le q$ and
\[
 q'\le F((1+a)q)\le (1+a)q-cq^2
\]
imply the same invariant upper bound used in Lemma~\ref{lem:spike-prefix}. Give state differences the weighted norm
\[
 \|(\Delta q,\Delta r)\|_a
       =|\Delta q|/\sqrt a+|\Delta r|.
\]
The preceding Jacobian estimates show that every step has operator norm at most $1+Ca$ in this norm. Consequently the whole prefix amplifies errors by at most
\[
 e^{CaH}\le e^{C(D+1)}.
\]
Changing between this norm and the ordinary coordinate norm costs at most $C\sqrt H$.

Now consider the padding and reverse-variance tail, whose schedule was chosen using the endpoint seed. Let $q_\ell^*$ denote its prescribed endpoint variance. The construction gives
\[
 q_\ell^*\ge ca,\qquad
 0\le\alpha_\ell\le Cq_\ell^*,
\]
and $q_\ell^*$ is constant during padding and nondecreasing during the reverse steps. At every other seed the variance satisfies
\[
 q_\ell(u)\le C_2 q_\ell^*.
\]
To verify this, the starting assertion follows from the prefix bounds. If $q\le C_2q_\ell^*$ then monotonicity and the quadratic upper drift give
\[
 F((1+\alpha_\ell)q)
 \le q+C\alpha_\ell q-cq^2.
\]
For a sufficiently large fixed $C_2$, the map points inward at $q=C_2q_\ell^*$. Since $q_{\ell+1}^*\ge q_\ell^*$, this proves induction. Choose the fixed amplitude cap small enough that all these envelopes stay in the small-variance interval.
Use the time-dependent norm
\[
 \|(\Delta q,\Delta r)\|_{q_\ell^*}
       =|\Delta q|/\sqrt{q_\ell^*}+|\Delta r|.
\]
One step has norm at most $1+Cq_\ell^*$; increasing the denominator at the next step can only decrease the norm. Padding has at most $C/a$ steps, and hence $\sum_{\rm padding}q_\ell^*=O(1)$. For a reverse step,
\[
 c(q_\ell^*)^2\le q_{\ell+1}^*-q_\ell^*
                         \le C(q_\ell^*)^2,
\]
so $\log(q_{\ell+1}^*/q_\ell^*)\ge cq_\ell^*$. Therefore
\[
 \sum_{\rm reverse}q_\ell^*\le C\log(1/a)\le C\log H.
\]
The tail's total amplification is at most a polynomial in $H$. This also covers the small-amplitude branch, which only pads at the endpoint variance and changes the readout.

Combining the stages, every inserted state error is amplified by at most
\[
 L_{D,H}\le H^C e^{CD}.
\]
The same estimates hold in enlarged envelopes for empirical states. They therefore imply the corresponding perturbation estimate by a standard first-exit induction: choose the accumulated error less than a fixed multiple of $a$, enlarge $C_1,C_2$ by a factor of two, and observe that the derivative bounds hold throughout every segment between an exact and an empirical state. The admissible state region $\{(q,r):r^2\le q\}$ is convex, so such segments remain admissible.

At each layer the conditional concentration estimate in Lemma~\ref{lem:spike-finite-width} gives coordinate error at most $C\epsilon$ simultaneously on all inputs, except with probability $Cm e^{-cN\epsilon^2}$. Choose
\[
 \epsilon=c\min\{a,\tau\}/(H L_{D,H}).
\]
The initial variance error can be required to obey the same bound, while its seed correlation is realized exactly. The initial Gram condition requires
$N\ge C\Lambda^{-2}\log(8m)+2m$, with $\Lambda=e^{-40(D+1)}$; this is $H^C e^{CD}$ after enlarging constants. The initial diagonal accuracy and the conditional layer failure estimates are achieved with total failure probability at most $c\tau^2$ once
\[
 N\ge C\epsilon^{-2}\log(C Hm/\tau).
\]
Both requirements follow from the claimed $C\tau^{-C}H^Ce^{CD}$ bound. On the good event the final hidden prediction error is at most a fixed multiple of $\tau$. On its complement the bounded readout factor and bounded correlations give a constant error bound. The native Gaussian readout contributes expected squared maximum at most $C\log(2m)/N$. Tightening the constants proves the stated expected supremum estimate.

Only the concentration widths have changed. The comparison law and its exact chain-rule entropy calculation are identical to Lemma~\ref{lem:spike-finite-width}, so the entropy bound is unchanged.
\end{proof}

\section{Completion of the finite-set comparison}
\begin{proof}[Proof of Theorem~\ref{thm:canonical-function-memorizer}]
Use the unit-sphere first-feature kernel with $\Lambda=e^{-40(D+1)}$, and seed the scalar endpoint with
\[
 q_1=F(2/(D+2)),\qquad r_1=\sqrt\Lambda/2.
\]
These obey the hypotheses of Lemma~\ref{lem:spike-prefix}, with fixed constants. Let $T=H-1$ be the number of internal transitions, and set
$\alpha=K(D+1)/T$ as in that lemma. Since $D\le\sqrt H/2$, its hypotheses hold for all sufficiently large $H$.
Reserve $L=\lceil C_1/\alpha\rceil$ final transitions, with a sufficiently large absolute $C_1$. Use the native burn and constant-spike prefix until transition $T-L$. Enlarge the fixed prefix constant $K$ so that $L\le T/4$; this is possible because $L/T=O(1/[K(D+1)])$. The prefix endpoint is therefore after $3T/4$. Its endpoint obeys
\[
 c\alpha\le q_0\le C\alpha,\qquad
 e_0=q_0-r_0^2\le q_0/2,
\]
and its cost is at most $C(D+1)^2/H$.

Write $a=\|f\|_\infty$. If $a\le2q_0$, hold the scalar variance exactly at $q_0$ for the remaining $L$ transitions, using
\[
 R^2-1=\frac{F^{-1}(q_0)-q_0}{q_0-e}.
\]
The padding lemma preserves $e\le\rho_1q_0$ for a fixed $\rho_1<1$, and the holding cost is $O(Lq_0^2)=O(\alpha)$. The final correlation has square at least $(1-\rho_1)q_0$, so the readout factor $B=a/r_H(1)$ has cost $B^2\le Ca^2/q_0\le Cq_0=O(\alpha)$.

If $a>2q_0$, use the sharp reverse-variance tail with target $Q=a/2$. Its variance sequence, and hence its length, depend only on $q_0,Q$; it uses at most $C_2/q_0\le C_2/(c\alpha)$ transitions. Choose $C_1$ larger than this constant. Pad the unused tail slots by holding at $q_0$ before the reverse path. The holding cost is again $O(\alpha)$ and preserves the defect invariant. The reverse-tail and readout bounds give total cost at most
\[
 4a+C_3a^2+C_3q_0+O(\alpha)
 \le4a+C_4A^2+C_4\alpha.
\]
In both cases all spike factors lie in $[1,2]$ after reducing $A_0$ and enlarging the depth threshold. The readout factor is bounded by an absolute constant. The final scalar correlation at seed value one is positive, even when $a=0$; in the latter case simply choose readout factor zero.

The schedule was selected using only its scalar endpoint. Its response to a general seed $u\in[-1,1]$ is continuous and odd. Corollary~\ref{cor:spike-ivt} therefore supplies seed values realizing every desired $f(x)$ exactly in the scalar recursion. Apply Theorem~\ref{thm:canonical-polynomial-width} to realize this schedule at the stated quantitative width with the required expected maximum squared error. Its exact entropy bound adds only $2\log2/N$ to the covariance and readout costs. Since
$\alpha=O((D+1)/H)$ is bounded by a constant times $(D+1)^2/H$, this proves~\eqref{eq:canonical-function-cost} and~\eqref{eq:canonical-function-approx}. For the lower bound, Corollary~\ref{cor:sup-entropy} with $2t=4-\delta$ gives the entropy bound in terms of $\E_\rho\|z\|_\infty$. The maximum-error assumption implies $\E_\rho\|z\|_\infty\ge\|f\|_\infty-\varepsilon$ by Cauchy--Schwarz, proving~\eqref{eq:canonical-function-lower}.

For the final assertion, choose $A$ so small that $C_0A^2\le\delta A/4$, then choose $c$ so small that $2C_0c^2\le\delta A/4$. With $D=\lfloor c\sqrt H\rfloor$, the remaining terms in $(D+1)^2/H$ beyond its limiting constant tend to zero. Increase $H$ and $N$ to absorb them and $2\log2/N$ into the remaining error allowance.
\end{proof}

\section{The same prior mechanism on an arbitrary fixed finite sphere subset}
\label{sec:arbitrary-finite-domain}
The explicit geometry of $\mathcal U_D$ is needed for a quantitative lower bound on the number of inputs. The maximum-amplitude mechanism itself is more general.

\begin{lemma}[Strict positivity of the native finite-set kernel]
\label{lem:finite-kernel-positive}
Let $\mathcal X$ be any nonempty finite set of distinct points on $\mathbb S^{d-1}$. The native first-layer kernel
\[
 K(x,x')=\E_A\tanh(A\cdot(1,x))\tanh(A\cdot(1,x')),
 \qquad A\sim\mathcal N(0,I/d),
\]
is strictly positive definite on $L^2(\mathcal X)$ for the uniform probability measure.
\end{lemma}
\begin{proof}
The Hermite calculation in Lemma~\ref{lem:sphere-direct-kernel}, with $\nu=0$, shows
\[
 K(x,x')=\sum_{j\text{ odd}}b_j^2
            \left(\frac{1+x\cdot x'}2\right)^j,
 \qquad b_j\ne0\quad\text{for every odd }j.
\]
Let $m=|\mathcal X|$, and choose an odd $j\ge m-1$. The tensor feature map for $(1+x\cdot x')^j$ contains, with positive coefficients, every coordinate monomial of total degree at most $j$. Choose a linear functional $\ell$ taking distinct values on the finitely many points of $\mathcal X$; the exceptional choices form a finite union of proper hyperplanes. Univariate Lagrange interpolation in $\ell(x)$ gives a polynomial of degree at most $m-1$ with any prescribed values on $\mathcal X$. Thus the degree-$j$ tensor features span every finite-set function, and their Gram matrix is positive definite. Its positive coefficient $b_j^2$ proves the claim.
\end{proof}

\begin{corollary}[An arbitrary finite domain with a quantitative kernel parameter]
\label{cor:arbitrary-finite-domain}
Let $\mathcal X\subset\mathbb S^{d-1}$ be finite, let $\Lambda>0$ be the minimum eigenvalue of its native kernel operator, and put
\[
 \Gamma=d+\log(2/\Lambda),\qquad G=\lceil\Gamma\rceil.
\]
There are absolute $A_0,C,H_0>0$ such that, for $H\ge H_0$, $G\le\sqrt H/2$, and every $f:\mathcal X\to[-A,A]$ with $0<A\le A_0$, every fixed $\varepsilon>0$ admits a comparison at all $N\ge C(1+\varepsilon^{-1})^C H^C e^{C\Gamma}$ satisfying
\[
 \E_\rho\max_{x\in\mathcal X}|z_w(x)-f(x)|^2\le\varepsilon^2,
\]
\[
 \frac2N\KL(\rho\|\mu)
 \le4\|f\|_\infty+CA^2+C\frac{\Gamma^2}{H}
       +\frac{2\log2}{N}.
\]
The finite-set entropy lower bound in Corollary~\ref{cor:sup-entropy} applies to the same domain. In particular, every fixed finite sphere subset eventually admits the same leading maximum-amplitude cost as depth and width grow.
\end{corollary}
\begin{proof}
Every first-layer preactivation has variance $2/d$, so its squared-output mean is $b_1=F(2/d)\ge c/d$ for an absolute $c>0$; the bound is uniform also for $d=1$. The kernel operator has trace $b_1\le1$, hence $\Lambda\le1/m$ and $\log m\le\log(1/\Lambda)\le\Gamma$. Lemma~\ref{lem:unit-seed} supplies a seed of size $r_1=\sqrt\Lambda/2$, with $r_1\ge e^{-C\Gamma}$ and $r_1\le\sqrt{b_1}/2$.

Apply Lemma~\ref{lem:spike-prefix} with the complexity parameter $G$ in place of $D$. Its hypotheses hold since $b_1\ge c/(G+1)$ and $\log(1/r_1)=O(G)$. The reverse-tail, padding, and intermediate-value arguments are identical to the proof of Theorem~\ref{thm:canonical-function-memorizer}. Finally $\log m+\log(1/\Lambda)\le2\Gamma=O(\sqrt H)$, so the weighted stability proof of Theorem~\ref{thm:canonical-polynomial-width}, with $G$ in place of $D$, gives the displayed polynomial-in-$H$ threshold with factor $e^{C\Gamma}$. Replacing $G$ by $\Gamma$ changes only absolute constants.
\end{proof}

\section{Actual function-neighborhood probabilities}
\label{sec:function-probabilities}
The comparison laws give bounds on the original prior's function probabilities, not only on a variational surrogate.

\begin{corollary}[Two-sided finite-set neighborhood costs]
\label{cor:function-neighborhood}
Use the hypotheses and notation of Theorem~\ref{thm:canonical-function-memorizer}, and put
\[
 a=\|f\|_\infty,\qquad
 \Delta=C_0A^2+C_0(D+1)^2/H,\qquad
 B_r(f)=\{w:\max_{x\in\mathcal U_D}|z_w(x)-f(x)|\le r\}.
\]
Fix $r>0$, $0<\eta<1$, and $0<\delta<1$. At every sufficiently large width $N\ge C_{r,\eta}H^{C_0}e^{C_0D}$,
\begin{align}
 &(4-\delta)(a-r)-C_\delta/H-2\log(2m)/N
 \ \le\ -\frac2N\log\mu(B_r(f))\notag\\
 &\hspace{38mm}\le
 \frac{4a+\Delta+4\log2/N}{1-\eta}.
 \label{eq:function-probabilities}
\end{align}
The lower estimate is valid at every width and does not require the comparison-width threshold.
\end{corollary}
\begin{proof}
On $B_r(f)$, the maximum output magnitude is at least $a-r$. The moment bound in the proof of Corollary~\ref{cor:sup-entropy} therefore gives
\[
 \mu(B_r(f))\le2m\exp\{ND_H(t)-tN(a-r)\},
 \qquad 2t=4-\delta.
\]
This proves the lower estimate, including when its right side is uninformative.

For the other direction, choose the comparison tolerance in Theorem~\ref{thm:canonical-function-memorizer} to be $r\sqrt\eta$. Markov gives $\rho(B_r(f))\ge1-\eta$. Binary relative entropy gives
\[
 \KL(\rho\|\mu)\ge(1-\eta)\log(1/\mu(B_r(f)))-\log2.
\]
Insert the comparison bound and multiply by $2/N$.
\end{proof}

At small maximum amplitude, the two estimates have the same leading cost after the relative neighborhood radius, comparison failure probability, and additive depth terms are made small. The corresponding unnormalized negative log probability is close to $2N\|f\|_\infty$. In this regime the leading bound is the same for simple and complicated functions on the specified finite input set.

\section{A fixed-accuracy learning obstruction at canonical mean field}
\label{sec:canonical-learning}
The matching supremum-norm bounds have a statistical consequence. They allow a deep canonical mean-field network to memorize arbitrary small labels almost as cheaply as it represents a smooth signal of the same amplitude. The resulting posterior is close to a simple capped empirical predictor, which preserves the signs of noisy singleton labels.

Write $\mathcal U_D$ for the unit-sphere input set used in Theorem~\ref{thm:canonical-function-memorizer}; it has $M=4\cdot2^D$ points and first coordinate $T\in\{-3/4,-1/4,1/4,3/4\}$. At canonical scaling, the first-layer, internal-layer, and readout variances are $1/d$, $1/N$, and $1/N^2$, respectively. For $n$ observations, define
\[
 \widehat R_S(z)=\frac1n\sum_{i=1}^n(z(X_i)-Y_i)^2,
 \qquad
 \Pi_{N,H,\lambda,S}(dw)=\frac{e^{-N\widehat R_S(z_w)/(2\lambda)}}{Z_{N,H,\lambda,S}}\,\mu_{N,H}(dw).
\]
Thus the Gaussian observation variance is $n\lambda/N$; every positive observation variance is allowed.

\begin{theorem}[Failure after exponentially many samples in square-root depth]
\label{thm:canonical-learning}
There are absolute constants $a,c,C>0$ and $H_0$ such that the following holds for every integer $H\ge H_0$. Set
\[
 D=\lfloor c\sqrt H\rfloor,\qquad d=D+2,\qquad n=4\cdot2^D,
\]
let $X$ be uniform on $\mathcal U_D$, and let
\[
 Y=g(X)+a\epsilon,\qquad
 g(X)=\frac a4(2+\tanh X_1),\qquad
 \Pr(\epsilon=1)=\Pr(\epsilon=-1)=\frac12,
\]
where $\epsilon$ is independent of $X$. With probability at least $9/10$ over $n$ IID observations, simultaneously for every integer $N\ge n^C$ and every $\lambda>0$, the ordinary posterior mean satisfies
\[
 R(\widehat f)\ge a^2+\frac{a^2}{4096}.
\]
The same bound holds for the mean of posterior draws clipped to $[-1,1]$, for clipping an ordinary posterior mean, and for arbitrary well-defined data-dependent probability mixtures of these predictions, including clipping after mixing.

The conditional mean $g$ has a width-two, one-hidden-layer tanh implementation with a fixed coefficient bound. It also has a bounded-width, bounded-coefficient implementation at every hidden depth. In particular, the excess-loss lower bound is a fixed positive number, while the sample count is $\exp(\Omega(\sqrt H))$.
\end{theorem}

The theorem concerns one fixed architecture at a time. Its event is uniform over widths and likelihood temperatures because it depends only on the occupancy and noise signs of the observed sample. No architecture prior, posterior computation assumption, or limiting interchange is involved.

\subsection{The proximal predictor selected by a supremum-norm penalty}
For a fixed sample, let $B_S\subset\mathbb R^n$ be the linear subspace of vectors constant on repeated copies of the same input. Equip it with $\|u\|_S^2=n^{-1}\sum_i u_i^2$. Let $b$ be the orthogonal projection of the label vector $y$ onto $B_S$; its entries are the empirical means at each observed input. Put $r_0=\|y-b\|_S^2$.

\begin{lemma}[A capped empirical predictor]
\label{lem:canonical-prox}
For every $\kappa>0$ and $\lambda>0$, the functional
\[
 J_\lambda(u)=r_0+\|u-b\|_S^2+\lambda\kappa\|u\|_\infty,
 \qquad u\in B_S,
\]
has a unique minimizer $u_\lambda$. There is a threshold $t_\lambda\in[0,\|b\|_\infty]$ for which
\[
 u_\lambda=\operatorname{clip}(b,[-t_\lambda,t_\lambda]).
\]
For all $u\in B_S$,
\[
 J_\lambda(u)-J_\lambda(u_\lambda)\ge\|u-u_\lambda\|_S^2.
\]
\end{lemma}
\begin{proof}
For a fixed cap $t\ge0$, the least-squares minimizer under $\|u\|_\infty\le t$ is coordinatewise clipping of $b$. Minimize $\|\operatorname{clip}(b,[-t,t])-b\|_S^2+\lambda\kappa t$ over $0\le t\le\|b\|_\infty$; no larger cap can improve the value. Strict convexity of the squared term gives uniqueness of $u_\lambda$. Finally $u\mapsto\lambda\kappa\|u\|_\infty$ is convex. At its minimizer, the first-order condition cancels the linear term in the expansion of the squared norm, leaving the displayed quadratic lower bound.
\end{proof}

\begin{lemma}[Posterior approximation by the capped predictor]
\label{lem:canonical-posterior-prox}
Suppose $|Y_i|\le2a$, and suppose the prior satisfies
\[
 \frac2N\operatorname{KL}(\rho\|\mu)
 \ge\kappa\,\mathbb E_\rho\|z\|_{\infty,\mathcal U_D}-e_0,
 \qquad 0<\kappa\le4.
\]
Suppose every function $v:\mathcal U_D\to[-2a,2a]$ has a comparison law with
\[
 \mathbb E\widehat R_S(z)\le\widehat R_S(v)+\zeta,
 \qquad
 \frac2N\operatorname{KL}(\rho_v\|\mu)\le4\|v\|_\infty+e_2.
\]
Then the Gaussian-likelihood posterior obeys
\[
 \mathbb E_\Pi\|z(S)-u_\lambda\|_S^2
 \le\zeta+\lambda\{2a(4-\kappa)+e_0+e_2\}.
 \tag{*}\label{eq:canonical-prox-gap}
\]
Here $u_\lambda$ is from Lemma~\ref{lem:canonical-prox}. The same upper bound holds for the squared distance between $u_\lambda$ and each of the following: the ordinary posterior mean, its clipping to $[-1,1]$, and the mean of draws clipped to that interval, provided $2a\le1$.
\end{lemma}
\begin{proof}
The Gibbs variational identity gives
\[
 F_S=\mathbb E_\Pi\widehat R_S(z)+\frac{2\lambda}{N}\operatorname{KL}(\Pi\|\mu)
 \le\mathbb E_{\rho_v}\widehat R_S(z)+\frac{2\lambda}{N}\operatorname{KL}(\rho_v\|\mu).
\]
Extend $u_\lambda$ by zero at unobserved atoms and use this extension as $v$. Since $\|u_\lambda\|_\infty\le\|b\|_\infty\le2a$, the comparison applies. The entropy lower bound and $\|z\|_{\infty,\mathcal U_D}\ge\|z(S)\|_\infty$ show
\[
 \mathbb E_\Pi J_\lambda(z(S))
 \le J_\lambda(u_\lambda)+\zeta+
       \lambda\{(4-\kappa)\|u_\lambda\|_\infty+e_0+e_2\}.
\]
Apply Lemma~\ref{lem:canonical-prox} pointwise and then integrate. Jensen proves the assertion for the mean. Projection onto $[-1,1]$ cannot increase distance to $u_\lambda\in[-2a,2a]$, proving both clipped assertions.
\end{proof}

\subsection{A sample witness common to every temperature and width}
\begin{proof}[Proof of Theorem~\ref{thm:canonical-learning}]
Fix $\Lambda=2^{14}$ and $\delta=2^{-44}$. Choose $t<2$ in Corollary~\ref{cor:sup-entropy} so close to two that $\kappa=2t\ge4-\delta$. First choose a fixed sufficiently small $a>0$, then $c>0$ sufficiently small. Theorem~\ref{thm:canonical-function-memorizer}, with amplitude bound $2a$, allows these choices so that, at $D=\lfloor c\sqrt H\rfloor$, every comparison function $v$ with $\|v\|_\infty\le2a$ has
\[
 \frac2N\operatorname{KL}(\rho_v\|\mu)\le4\|v\|_\infty+e_2,
 \qquad e_2\le\delta a/2,
\]
for all sufficiently large $H$ and all $N\ge n^C$. Use the fixed RMS approximation tolerance $\varepsilon=2^{-28}a$. Since $|v-Y|\le4a$, the corresponding empirical-risk error satisfies
\[
 \zeta\le8a\varepsilon+\varepsilon^2.
\]
The finite-set entropy floor has
\[
 e_0=2D_H(t)+\frac{2\log(2n)}N\le\delta a/2
\]
after increasing the depth threshold and the width constant. Increase them again so that $1/N\le a^2$. The comparison-width requirement is a fixed polynomial in $H$ and $e^D$ at this fixed accuracy. Since $n=4\cdot2^D$ and $D=\lfloor c\sqrt H\rfloor$, it is bounded by $n^C$ after increasing the absolute exponent and depth threshold. All these choices are fixed constants, independent of the observed sample.

For $0<\lambda\le\Lambda a$, Lemma~\ref{lem:canonical-posterior-prox} now gives
\[
 \mathbb E_\Pi\|z(S)-u_\lambda\|_S^2
 \le8a\varepsilon+\varepsilon^2+3\Lambda\delta a^2
 <10^{-6}a^2.
 \label{eq:canonical-small-temp}
\]
Thus all three stated posterior predictions lie within $a/1000$ of $u_\lambda$ in empirical norm.

For $\lambda>\Lambda a$, use the prior itself as the variational comparison. Its output is centered and has variance at most $1/N$ at every input. Since $|Y|<2a$,
\[
 F_S\le\widehat R_S(0)+1/N\le5a^2.
\]
Dropping the nonnegative empirical risk from the entropy lower bound yields
\[
 \mathbb E_\Pi\|z\|_{\infty,\mathcal U_D}
 \le\frac{5a^2}{\kappa\lambda}+\frac{e_0}{\kappa}
 <\frac a{1000}.
\]
The same bound applies to the sup norm of each stated posterior prediction. Hence large temperatures give predictors uniformly close to zero, itself a capped predictor with threshold zero.

Let $A_S$ be the set of sample indices whose input occurs exactly once and whose independent noise sign is negative. Since $n$ observations are drawn from $n$ equally likely atoms, the expected singleton count is $n(1-1/n)^{n-1}$. Changing one input alters this count by at most two, so bounded differences shows that it is at least $n/3$ with probability tending to one. Conditional on the inputs, the signs on these singleton indices are independent fair signs. A Chernoff bound therefore gives
\[
 |A_S|\ge n/10
\]
with probability at least $9/10$ for all sufficiently large $H$. Work on this event. Define $w_i=-\mathbf1_{i\in A_S}$. At such an index the empirical mean is $b_i=Y_i=g(X_i)-a<0$. Every capped predictor therefore has $u_{\lambda,i}\le0$, whereas $g(X_i)>a/4$. Consequently
\[
 \langle w,u_\lambda-g\rangle_S\ge\frac a4\frac{|A_S|}n\ge\frac a{40}.
\]
The small-temperature approximation and $\|w\|_S\le1$ show that each posterior prediction $f$ satisfies
\[
 \langle w,f-g\rangle_S\ge\frac a{40}-\frac a{1000}>\frac a{64}.
 \label{eq:canonical-witness}
\]
The same follows directly from the uniform near-zero bound at large temperatures.

This linear witness is common to every admissible width and temperature. It therefore persists under arbitrary well-defined probability mixtures. More precisely, each small-temperature prediction admits a decomposition $f=u+e$, where $u\le0$ on $A_S$ and $\|e\|_S\le a/1000$; large-temperature predictions use $u=0$. Mixtures preserve these two properties by convexity. Since clipping is monotone and fixes the interval containing every capped predictor, clipping after mixing also preserves the same approximation and witness. For a general mixture, the reference capped predictors remain in $[-2a,2a]$; Jensen and the contraction of interval projection justify this last claim.

Let $h(x)=-1$ on the distinct singleton atoms represented by $A_S$ and zero elsewhere. The population law assigns each such atom mass $1/n$, so
\[
 \langle h,f-g\rangle_{L^2(P_X)}=\langle w,f-g\rangle_S>\frac a{64},
 \qquad \|h\|_{L^2(P_X)}^2=|A_S|/n\le1.
\]
Cauchy--Schwarz gives $\|f-g\|_{L^2(P_X)}^2\ge a^2/4096$. Independent centered label noise gives $R(f)=a^2+\|f-g\|_{L^2(P_X)}^2$, proving the risk bound.

Finally the width-two shallow representation uses one constant tanh feature and one feature $\tanh T$. To obtain a fixed-width representation at arbitrary depth, preserve the constant feature at a nonzero tanh fixed point and preserve $(\tanh T,\tanh(2T))$ on the four possible values of $T$. The two feature values determine $T$ by a fixed linear map: the ratio $\tanh(2T)/\tanh T$ is strictly different at $|T|=1/4$ and $3/4$. Applying that map and then the two tanh features reproduces the same pair at every layer. All resulting coefficients are fixed constants. This completes the proof.
\end{proof}

The conclusion is specific about what must fail. It rules out a bound polynomial jointly in learner depth and the complexity of an approximating teacher for this unrestricted canonical Gaussian learner. It is compatible with polynomial sample complexity at every fixed depth, and it does not assert an exponential lower bound in teacher depth when the learner itself remains shallow. The independent label noise is essential to this construction; the analogous noiseless Boolean question requires a different argument.

\section{The obstruction with Gaussian label noise}
\label{sec:canonical-Gaussian-noise}
The same failure occurs when the conditional label distribution is Gaussian, so that the observation likelihood has the correct functional form. The likelihood variance may still be chosen arbitrarily, including its true value.

\begin{corollary}[A correctly specified Gaussian observation model]
\label{cor:canonical-Gaussian-noise}
There are absolute $a,c,C>0$ and $H_0$ such that the construction of Theorem~\ref{thm:canonical-learning} with
\[
 Y=g(X)+aZ,\qquad Z\sim\mathcal N(0,1)\text{ independent of }X,
 \qquad g(X)=\frac a4(2+\tanh X_1),
\]
satisfies, with probability at least $9/10$,
\[
 R(\widehat f)\ge a^2+\frac{a^2}{512^2}
\]
simultaneously for every $N\ge n^C$ and every positive likelihood variance, after $n=4\cdot2^{\lfloor c\sqrt H\rfloor}$ IID samples. The clipped predictions and arbitrary mixtures listed in Theorem~\ref{thm:canonical-learning} satisfy the same bound.
\end{corollary}
\begin{proof}
We give the changes needed to accommodate unbounded labels. Set
\[
 A=16a,\qquad \Lambda_0=2^{20},\qquad
 \delta=2^{-80},\qquad \varepsilon=2^{-50}a.
\]
Choose the fixed amplitude $a$ sufficiently small and then $c$ sufficiently small. Theorem~\ref{thm:canonical-function-memorizer}, now with comparison amplitude bound $A$, and Corollary~\ref{cor:sup-entropy} ensure, at all sufficiently large $H$ and all $N\ge n^C$,
\[
 \kappa\ge4-\delta,\qquad e_0\le\delta a/2,
 \qquad e_2\le\delta a/2,
\]
with the notation of Lemma~\ref{lem:canonical-posterior-prox}. Every comparison output bounded by $A$ has expected maximum squared approximation error at most $\varepsilon^2$ and entropy cost at most $4\|v\|_\infty+e_2$. The fixed-accuracy comparison threshold is polynomial in $H$ and $e^D$, and is therefore covered by $N\ge n^C$ for an increased absolute exponent and sufficiently large $H$. Increase the thresholds so that $1/N\le a^2$.

For any observed labels, let $b$ and $u_\lambda$ be as in Lemma~\ref{lem:canonical-prox}. The minimizing cap may now exceed $A$. Define
\[
 v_\lambda=\operatorname{clip}(u_\lambda,[-A,A]),\qquad
 T_S=\frac1n\sum_i(|b_i|-A)_+^2.
\]
Since $u_\lambda$ is itself a clipped version of $b$, reducing its cap gives
\begin{equation}
 J_\lambda(v_\lambda)-J_\lambda(u_\lambda)\le T_S.
 \label{eq:Gaussian-cap-error}
\end{equation}
Indeed the supremum penalty decreases. At an atom where clipping changes the value, the increase in its squared-error term is
$(|b|-A)^2-(|b|-t_\lambda)_+^2\le(|b|-A)^2$.
The function $y\mapsto(|y|-A)_+^2$ is convex, so averaging within repeated-input groups also gives
\[
 T_S\le\frac1n\sum_i(|Y_i|-A)_+^2.
\]

Work first on the event $\widehat R_S(0)\le6a^2$. A comparison law for $v_\lambda$, extended by zero at unseen inputs, changes its empirical risk by at most
\[
 \zeta\le2\|v_\lambda-Y\|_S\varepsilon+\varepsilon^2
 \le2(16+\sqrt6)a\varepsilon+\varepsilon^2
 <38a\varepsilon+\varepsilon^2.
\]
Repeating the variational argument in Lemma~\ref{lem:canonical-posterior-prox}, using~\eqref{eq:Gaussian-cap-error}, yields
\[
 \E_\Pi\|z(S)-u_\lambda\|_S^2
 \le T_S+\zeta+\lambda\{(4-\kappa)A+e_0+e_2\}.
\]
Consequently, on the further event $T_S\le2^{-50}a^2$, every
$0<\lambda\le\Lambda_0a$ obeys
\[
 \E_\Pi\|z(S)-u_\lambda\|_S^2
 \le2^{-50}a^2+38\cdot2^{-50}a^2+2^{-100}a^2
            +17\cdot2^{-60}a^2
 <2^{-40}a^2.
\]
The ordinary posterior mean is therefore within $a/10000$ of $u_\lambda$ in empirical norm. Each clipped prediction is within the same distance of $\operatorname{clip}(u_\lambda,[-1,1])$, by contraction and Jensen. These reference predictions are nonpositive wherever a singleton label is negative.

For $\lambda>\Lambda_0a$, the prior comparison gives $F_S\le7a^2$. The entropy floor then implies
\[
 \E_\Pi\|z\|_{\infty,\mathcal U_D}
 \le\frac{7a^2}{\kappa\lambda}+\frac{e_0}{\kappa}
 <a/10000.
\]
All the stated posterior predictions are uniformly within this distance of zero.

It remains to check that the sample events have high probability. Since $|g|\le a$, the event $n^{-1}\sum_iZ_i^2\le2$ implies $\widehat R_S(0)\le6a^2$ and has probability tending to one. Also
\[
 \E\frac{T_S}{a^2}
 \le\E(|Z|-15)_+^2
 \le\E[Z^2\mathbf1_{|Z|>15}]
 =2\{15\varphi(15)+\overline\Phi(15)\}
 <2^{-150},
\]
where $\varphi$ and $\overline\Phi$ are the standard Gaussian density and upper tail; the last bound follows from $\overline\Phi(t)\le\varphi(t)/t$. Markov gives
$\Pr(T_S>2^{-50}a^2)<2^{-100}$.

With probability tending to one, at least $n/3$ observations are singleton inputs. On these observations the noises are independent, and
\[
 \Pr(-2\le Z\le-1)\ge e^{-2}/\sqrt{2\pi}>1/20.
\]
A Chernoff bound therefore ensures at least $n/100$ singleton observations with $Z\le-1$, with probability tending to one. Let their index set be $A_S$. At each such observation, $Y=g+aZ<0$, whereas $g>a/4$. All four events hold jointly with probability at least $9/10$ for sufficiently large $H$.

Use the witness $w_i=-\mathbf1_{i\in A_S}$. Every capped empirical reference, and its clipping, has nonpositive coordinates on $A_S$, so every stated posterior prediction satisfies
\[
 \langle w,\widehat f-g\rangle_S
 \ge a/400-a/10000>a/512.
\]
The same witness persists under arbitrary probability mixtures and clipping after mixing: their reference predictions remain nonpositive on $A_S$, and their empirical approximation errors remain at most $a/10000$ by convexity and contraction. The references are bounded on the finite observed sample by its largest label magnitude, so this argument also covers well-defined general mixtures.

Each singleton atom has population probability $1/n$. Translating this witness to the population and applying Cauchy--Schwarz as in Theorem~\ref{thm:canonical-learning} gives
$\|\widehat f-g\|_{L^2(P_X)}^2\ge a^2/512^2$. Finally $R(\widehat f)=a^2+\|\widehat f-g\|_2^2$ for the independent Gaussian noise.
\end{proof}

At the true observation variance $a^2$, the corresponding likelihood parameter is $\lambda=Na^2/n$. The corollary allows this choice as well as every other positive value, including data-dependent choices. Thus the obstruction is not caused by a mismatch between the likelihood family and the conditional noise distribution.

% Research module: additive depth-baseline and amplitude entropy floor.
\section{A joint depth and amplitude cost below criticality}
\label{sec:joint-prefactor-entropy}
A polynomially weighted change of measure compares the entire subcritical variance chain with the canonical chain. It retains the depth cost while also retaining the sharp amplitude penalty.

\begin{lemma}[A variance-chain kernel domination]
\label{lem:joint-kernel-domination}
Let $N\ge2$, $0<\nu\le1$, and put
\[
 \beta=1-\nu,\qquad p=\frac{N\beta}{2\nu},\qquad
 \kappa(\nu)=1-\nu+\nu\log\nu,\qquad
 M_\nu=\exp\{-N\kappa(\nu)/(2\nu)\}.
\]
For the native scalar variance transition
\[
 q'=\frac1N\sum_{i=1}^N\tanh^2(\sqrt{\nu q}\,G_i)
\]
and every nonnegative measurable function $f$,
\[
 \E[(q')^p f(q')\mid q]
 \le M_\nu q^p\E[f(\widetilde q)\mid q],
 \qquad
 \widetilde q=\frac1N\sum_{i=1}^N\tanh^2(\sqrt q\,G_i).
\]
Thus the polynomially weighted subcritical transition, normalized by $M_\nu$, is dominated as a positive kernel by the canonical transition.
\end{lemma}
\begin{proof}
For $q>0$ let $S=N^{-1}\sum_iG_i^2$. Then $q'\le\nu qS$. Concavity of the logarithm, with tangent at $S=1/\nu$, gives
\[
 (\nu S)^p\le
 \exp\{p\nu(S-1/\nu)\}
 =\exp\left\{\frac\beta2\sum_iG_i^2-\frac{N\beta}{2\nu}\right\}.
\]
Multiplication by the standard Gaussian density changes its quadratic exponent to $-\nu\sum_iG_i^2/2$. The normalization factor is
\[
 \nu^{-N/2}e^{-N\beta/(2\nu)}
       =e^{-N\kappa(\nu)/(2\nu)}=M_\nu.
\]
Under the resulting product law, $G_i$ has variance $1/\nu$, so $\sqrt{\nu q}G_i$ becomes a canonical $N(0,q)$ preactivation. Multiplying by $q^p$ proves the inequality. The cases $q=0$ or $p=0$ follow directly, with the usual convention $q^0=1$.
\end{proof}

\begin{corollary}[A joint fractional and exponential moment]
\label{cor:joint-variance-moment}
For $3/2\le\theta<2$, let $k_\theta>0$ be the constant in the canonical heat supersolution. Then, for any native first layer with $q_1\le1$,
\[
 \E[q_H^p e^{\theta Nq_H}]
 \le\exp\left\{-\frac{N\kappa(\nu)(H-1)}{2\nu}
             +\frac{\theta N}{1+k_\theta(H-1)}\right\}.
\]
\end{corollary}
\begin{proof}
For $q>0$, divide the preceding lemma by $M_\nu q^p$ to obtain a subprobability transition kernel dominated by the canonical transition. Iteration preserves domination for every nonnegative test function. The polynomial weights telescope along a path, giving $M_\nu^{H-1}q_1^p$ times a canonical-chain expectation started at $q_1$. The canonical heat-supersolution estimate bounds its exponential moment uniformly over $q_1\le1$ by $\exp\{\theta N/[1+k_\theta(H-1)]\}$. Finally $q_1^p\le1$; integrate over the first layer.
\end{proof}

\begin{theorem}[An additive baseline-plus-amplitude entropy floor]
\label{thm:joint-prefactor-floor}
Let $N\ge2$, let $\mathcal X$ be any finite input set of size $m\ge1$, put $Z_*=\max_{x\in\mathcal X}|z_w(x)|$, and use the Gaussian tanh prior with $s=\nu/N$, $0<\nu\le1$. Fix $0<u<1$ and $0<t<2$. Choose $3/2\le\theta<2$ with $t^2<2\theta$, and write
\[
 L_{t,\theta}=\log\frac1{1-t^2/(2\theta)},\qquad
 B=\kappa(\nu)(H-1).
\]
Every comparison law $\rho$ satisfies
\begin{align*}
 2s\operatorname{KL}(\rho\|\mu)
 &\ge B\,\rho\{Z_*>u\}+2\nu t\,\E_\rho Z_* -\mathcal E,\\
 \mathcal E
 &=\frac{2\nu\theta}{1+k_\theta(H-1)}
   +\left(\beta+\frac\nu N\right)L_{t,\theta}
   +2\beta\log(1/u)
   +\frac{2\nu\log(2m)}N+2\nu t u.
\end{align*}
The expectation can be interpreted by monotone truncation. In particular, at $\beta=O(H^{-1/2})$ and fixed $u,t,\theta$, all error terms other than $2\nu t u$ vanish as $H,N\to\infty$ with $\log m=o(N)$.
\end{theorem}
\begin{proof}
At a fixed input the readout is conditionally $z=\sqrt{\nu q_H/N}\,G$, with an independent standard Gaussian $G$. Young's inequality gives
\[
 tN|z|\le\theta Nq_H+\frac{t^2\nu}{4\theta}G^2.
\]
Set $r=t^2\nu/(2\theta)<1$. Gaussian integration yields
\[
 \E|G|^{2p}e^{rG^2/2}
 =(1-r)^{-p-1/2}\E|G|^{2p}.
\]
Also $(\E|G|^{2p})^{1/(2p)}\le\sqrt{2p+2}$ for $p>0$ by the even-integer moment formula and monotonicity of moments. Since
\[
 (\nu/N)(2p+2)=\beta+2\nu/N\le1,
\]
we obtain
\[
 \E\bigl[|z|^{2p}e^{tN|z|}\bigr]
 \le(1-r)^{-p-1/2}\E[q_H^p e^{\theta Nq_H}].
\]
This inequality also holds when $p=0$ directly. On $|z|>u$, divide by $u^{2p}$. Corollary~\ref{cor:joint-variance-moment} and $1-r\ge1-t^2/(2\theta)$ then give
\[
 \E[e^{tN|z|}\mathbf1_{|z|>u}]
 \le\exp\left\{-\frac{NB}{2\nu}+\frac{N e}{2\nu}\right\},
\]
where
\[
 e=\frac{2\nu\theta}{1+k_\theta(H-1)}
    +\left(\beta+\frac\nu N\right)L_{t,\theta}
    +2\beta\log(1/u).
\]
The exponential of a maximum on its threshold event is at most the sum of the corresponding thresholded exponentials. Hence the same bound for $Z_*$ has $e$ replaced by $e+2\nu\log m/N$.

Let $A=\{Z_*>u\}$ and use the entropy test potential
\[
 V=\left(\frac{NB}{2\nu}+tNZ_*\right)\mathbf1_A.
\]
Its prior exponential moment is at most
$1+\exp\{N[e+2\nu\log m/N]/(2\nu)\}$, so its logarithm is at most $N[e+2\nu\log m/N]/(2\nu)+\log2$. The entropy variational inequality, multiplied by $2s=2\nu/N$, gives
\[
 2s\operatorname{KL}\ge B\rho(A)+2\nu t\E_\rho[Z_*\mathbf1_A]
                -e-2\nu\log(2m)/N.
\]
Finally $\E[Z_*\mathbf1_A]\ge\E Z_*-u$. This proves the theorem. For infinite moments, truncate the test potential and then pass to the limit.
\end{proof}

% Research module: a uniform function comparator for every 0 < Ns <= 1.
\section{A uniform function comparison across all subcritical prefactors}
\label{sec:all-prefactor-function-comparison}
Inflating only the first-layer covariance removes the degeneration of the native feature kernel at very small prefactors. The remaining construction pays the same depth baseline as the joint entropy floor.

\begin{theorem}[A baseline-plus-amplitude finite-function comparison]
\label{thm:all-prefactor-function-comparison}
There are absolute constants $A_0,c,C>0$ such that the following holds. Let $H$ be sufficiently large, $D\ge1$, $d=D+2$, and
\[
 E=D+\log H\le cH,
 \qquad 0<\nu\le1,\quad \beta=1-\nu,\quad s=\nu/N.
\]
For every function $f:\mathcal U_D\to[-A,A]$ with $0<A\le A_0$, every tolerance $0<\tau<1$, and every width
\[
 N\ge C\tau^{-C}H^C e^{CD},
\]
there is a comparison law satisfying
\begin{align*}
 2s\operatorname{KL}(\rho\|\mu_s)
 &\le \kappa(\nu)(H+d)+4\|f\|_\infty+CA^2
           +C\left\{\beta E+\frac{E^2}{H}\right\},\\
 \left(\E_\rho\max_{x\in\mathcal U_D}|z_w(x)-f(x)|^2\right)^{1/2}
 &\le\tau,
\end{align*}
where $\kappa(\nu)=1-\nu+\nu\log\nu$. All constants are uniform over the entire interval $0<\nu\le1$.
\end{theorem}
\begin{proof}
\emph{First-layer encoding.}
Draw the first-layer entries with variance $1/d$, rather than their native variance $\nu/d$. Their normalized relative entropy is exactly
\[
 \frac{2\nu}{N}\frac{N(d+1)}2
       \{\nu^{-1}-1-\log(\nu^{-1})\}
       =(d+1)\kappa(\nu).
\]
The Gaussian feature kernel on $\mathcal U_D$ has minimum eigenvalue at least $\Lambda=e^{-40(D+1)}$, and its common diagonal is $q_1=F(2/d)\asymp1/(D+1)$. Let $\zeta_D=\sqrt\Lambda/2$. With probability arbitrarily close to one at the stated width, the empirical kernel is at least half this kernel and the diagonal is uniformly close to $q_1$. On this good event choose a measurable unit direction realizing
\[
 e_1^{\mathsf T}h^1(x)/\sqrt N=\zeta_Du(x)
\]
for each prescribed seed $u(x)\in[-1,1]$. The minimum-norm/nullspace construction proves this exactly. On the bad event choose an arbitrary fixed unit direction. We do not condition the first-layer distribution. Consequently its entropy is exactly the displayed covariance-inflation cost, while the bad event is handled in the final expected-error estimate.

\emph{A controlled burn and growing prefix.}
At an internal layer write $W=\nu R^2$ for the effective directional variance. The scalar recursion is
\[
 v=\nu q+(W-\nu)r^2,\qquad q'=F(v),\qquad
 r'=\sqrt W\,r(1-q').
\]
First use $W=1$, equivalently $R^2=1/\nu$, for $\lfloor H/4\rfloor$ transitions. Since $v\le q$, the variance is bounded above by the canonical native recursion. Hence $q\le C/H$ after this burn and $\sum q\le C\log H$ during it. Also $r'=r(1-q')$, so the endpoint seed correlation remains at least $\exp[-C(D+\log H)]$, independently of $\nu$. In particular, the logarithmic variance and coherence deficits are both $O(E)$, since $q\ge r^2$ and $q\le1$.

Put $a=KE/H$ for a sufficiently large absolute $K$, and choose $c$ so that $a$ is small. During the growing stage use $W=1+a$. Then
\[
 v=\nu q+(a+\beta)r^2\le(1+a)q,
\]
which gives the invariant upper bound $q\le C_1a$. For $\theta=r^2/q$,
\[
 \theta'=\frac{(1+a)\theta}{\nu+(a+\beta)\theta}
                  \frac{v(1-F(v))^2}{F(v)}.
\]
The last factor is between $1-Ca^2$ and one. While $\theta\le1/2$, the first factor improves it by at least $1+ca$; therefore it reaches one half within $CE/a$ steps. At that point it stays bounded away from zero, and its deficit obeys
\[
 1-\theta'\le[1-c(a+\beta)](1-\theta)+Ca^2.
\]
Thus, after at most $C\log H/a$ more steps, $1-\theta\le Ca$, and also
\[
 1-\theta\le \frac{a}{4(a+\beta)}.
\]
For the latter bound, the stationary error $Ca^2/(a+\beta)$ is smaller than the requested threshold at small $a$; the relaxation time is at most $C/a$, since $\log(1+\beta/a)/(a+\beta)\le C/a$. The two bounds can be required simultaneously by taking the larger of their relaxation times.

With this coherence, $v/q=1+a-(a+\beta)(1-\theta)\ge1+3a/4$. Whenever $q$ is below a sufficiently small multiple of $a$, the quadratic lower drift of $F$ gives $q'\ge(1+ca)q$. The correlation cannot have lost more than $CE$ in its logarithm during the preceding growing steps: the coherence stages just established use only $CE/a$ transitions, and each loses at most $Ca$ in the logarithm. The constant in this $O(E)$ loss is independent of the later choice of $K$. Hence $q\ge e^{-CE}$ throughout those steps, and it reaches $ca$ in another $CE/a$ transitions. Increasing $K$ makes all these stages finish before time $3H/4$. The bounds
\[
 ca\le q\le Ca,\qquad e:=q-r^2\le Ca^2\le C'q^2
\]
then remain valid. Continue the growing stage until $H-1-\lceil C_0/a\rceil$ internal transitions have been allocated, reserving the last $\lceil C_0/a\rceil$ slots for padding and reverse transport. This time is after $3H/4$ at sufficiently large $H$.

The normalized action of a layer with effective variance $W$ is
\[
 \mathcal A_\nu(W)=W-\nu-\nu\log(W/\nu).
\]
It equals $\kappa(\nu)$ at $W=1$, and exactly
\[
 \mathcal A_\nu(1+a)-\kappa(\nu)
       =a-\nu\log(1+a)\le\beta a+a^2/2.
\]
Thus the prefix pays its baseline $\kappa(\nu)$ per transition plus at most $C\{\beta E+E^2/H\}$.

\emph{Reverse transport and readout.}
The exact defect identity is
\[
 e'=J(q')+\nu e(1-q')^2,
 \qquad J(q)=\tfrac23q^3+O(q^4),\quad J\ge0.
\]
The invariant $e\le K_0q^2$ is preserved for a sufficiently large fixed $K_0$. Indeed, for a reverse step $q=F(q')$, the small-variance expansion gives $(q')^2-q^2(1-q')^2\ge3(q')^3$, whereas $J(q')\le2(q')^3$. Thus $e'\le K_0(q')^2$ for $K_0\ge2$; the factor $\nu\le1$ only helps. At a padding step $q'=q$, the same conclusion follows from $q^2[1-(1-q)^2]\ge(3/2)q^3$ after increasing $K_0$ if necessary. For a reverse step prescribe $q'=F^{-1}(q)$ and use
\[
 W=\nu R^2=\frac{F^{-1}(q')-\nu e}{q-e}.
\]
Uniformly over $\nu\in(0,1]$,
\[
 W-1=4q+O(q^2+\beta q),\qquad q'-q=2q^2+O(q^3).
\]
Expanding $\mathcal A_\nu(W)$ at $W=1$ therefore gives the upper bound
\[
 \mathcal A_\nu(W)
 \le\kappa(\nu)+4(q'-q)+Cq(q'-q)+C\beta q.
\]
There are at most $C/q_0$ reverse steps from $q_0\asymp a$, and their variances satisfy $\sum q\le C\log H$. Their action above baseline is at most
$4(Q-q_0)+CQ^2+C\beta\log H$.
Padding at $q_0$ has $W-1=O(q_0)$ and costs at most $C(\beta q_0+q_0^2)$ above baseline per step. The reserved number of slots is $O(1/a)$, so padding adds only $C(\beta+a)$ above baseline.

For $b=\|f\|_\infty>2q_0$, choose the reverse endpoint $Q=b/2$. The final variance is $Q+O(Q^2)$ and defect $O(Q^2)$, so readout factor $B=b/r_H(1)$ costs $B^2=2b+O(b^2)$. Reverse action above baseline plus readout cost is at most $4b+CA^2+C\beta\log H+Ca$. When $b\le2q_0$, use only holding and the readout $B=b/r_H(1)$; its additional cost is $O(\beta+a)$. All $O(a)$ terms are absorbed into $CE^2/H$, and all $O(\beta)$ terms into $C\beta E$. The scalar endpoint is continuous and odd in its seed, so the intermediate value theorem supplies seeds for all values of $f$.

There are $H-1$ internal transitions. Adding their baseline to the first-layer cost gives exactly $\kappa(\nu)(H+d)$. The readout factor is bounded by an absolute constant on both branches.

\emph{Uniform finite-width error.}
At each internal layer draw rows with covariance
\[
 \frac1N[\nu I+(W_\ell-\nu)e_{\ell-1}e_{\ell-1}^{\mathsf T}],
\]
and define the new unit direction using its standardized Gaussian directional coordinates. All effective variances $W_\ell$ are bounded by an absolute constant, even when the relative inflation $R_\ell^2=W_\ell/\nu$ is large. Conditional concentration for the variance, correlation, and direction normalization therefore has constants uniform down to $\nu\downarrow0$. The readout mean is $Be_H/\sqrt N$ and its covariance is the native $\nu I/N^2$. Its entropy cost is exactly $B^2$ after normalization, and its expected squared maximum noise is $O(\nu\log m/N)$.

For clarity, the scalar propagation estimate can also be uniform in $\nu$. In a stage with envelope $q\le Cb$ and $W=1+O(b)$, use
\[
 \|(\Delta q,\Delta r)\|_{b,\beta}
 =w_{b,\beta}|\Delta q|+|\Delta r|,
 \qquad w_{b,\beta}=\frac{K_1\sqrt b}{\beta+b}.
\]
Here $K_1$ is a sufficiently large fixed constant. The Jacobian entries have absolute bounds
\[
 \begin{pmatrix}
 \nu & C(\beta+b)\sqrt b\\
 C\sqrt b & 1+Cb
 \end{pmatrix}.
\]
The weighted first column sum is at most
$\nu+(C/K_1)(\beta+b)\le1+Cb$, while the second is at most $1+CK_1b$. Thus one step has weighted norm at most $1+Cb$.
During the burn take $b_\ell=C/(1+c\ell)$, which bounds every seed variance. The total $\sum b_\ell$ is $O(\log H)$. The weight changes have total logarithmic variation $O(\log H)$, since $b_\ell$ is monotone and $w_{b,\beta}=K_1\sqrt b/(\beta+b)$. During growth take $b=a$, giving amplification at most $e^{CaH}=e^{CE}$. During the tail take $b=q_\ell^*$, the endpoint variance; every other seed satisfies $q_\ell\le Cq_\ell^*$ by quadratic drift, and $\sum q_\ell^*=O(\log H)$. Its weights again have total logarithmic variation $O(\log H)$. The finitely many initial steps outside the small-variance interval change only constants. Conversion to ordinary coordinate norms costs at most a polynomial in $H$.
Consequently every state error is amplified by at most $H^Ce^{CD}$, uniformly in $\nu$. Enlarged envelopes and a first-exit induction justify the same estimate for empirical states. Conditional Gaussian concentration, the first-layer kernel concentration, and a union bound now give the stated width $C\tau^{-C}H^Ce^{CD}$ and expected supremum error. The bad first-layer event contributes only a bounded error times its arbitrarily small probability.

Finally, the chain rule adds the first-layer covariance entropy, the deterministic action $\mathcal A_\nu(W_\ell)$ at each internal layer, and $B^2$ for the readout. The directions may depend measurably on all preceding parameters, but their covariance eigenvalues are prescribed, so the action is exact and there are no conditioning terms. This proves the theorem.
\end{proof}

\section{A baseline cost plus an amplitude penalty}
\label{sec:baseline-prox}
The statistical argument also permits a positive cost for producing any appreciable output. Its value may vary with the prior scale. The conclusion needed for the noise witness is unchanged: the posterior mean is close to a mixture of zero and a capped empirical predictor.

\begin{lemma}[A baseline-plus-amplitude posterior bound]
\label{lem:baseline-prox}
Let the input support be a finite set $\mathcal U$, let $|Y_i|\le2a$, and write $M(z)=\max_{x\in\mathcal U}|z(x)|$. Suppose $B\ge0$, $0<c\le4$, and $t>0$ satisfy the entropy bound
\[
 \mathfrak s\operatorname{KL}(\rho\|\mu)
 \ge B\Pr_\rho(M(z)>t)+c\,\mathbb E_\rho M(z)-e_0
\]
for every absolutely continuous comparison law, where $\mathfrak s>0$ is the entropy normalization and $s,\zeta,e_0,e_2\ge0$. Suppose the prior itself has expected empirical loss at most $\widehat R(0)+s$, and every function $v:\mathcal U\to[-2a,2a]$ has a comparison satisfying
\[
 \mathbb E\widehat R(z)\le\widehat R(v)+\zeta,
 \qquad \mathfrak s\operatorname{KL}(\rho_v\|\mu)
        \le B+4\|v\|_\infty+e_2.
\]
Let $\Pi$ be the Gibbs posterior minimizing
$\mathbb E\widehat R+\lambda\mathfrak s\operatorname{KL}$, and let $u_\lambda$ minimize
\[
 J_\lambda(u)=\widehat R(u)+\lambda c\|u\|_\infty
\]
on sample vectors constant on repeated inputs. Then $u_\lambda$ is a capped empirical-mean predictor. For some $p\in[0,1]$,
\[
 \|\mathbb E_\Pi z(S)-p u_\lambda\|_S
 \le t+\sqrt{4at+s+\zeta+
          \lambda\{2a(4-c)+e_0+e_2\}}.
\]
If $2a\le1$, the same bound holds for the mean of draws clipped to $[-1,1]$ and for clipping the raw posterior mean. The reference vector $p u_\lambda$ is nonpositive at every observed input with a nonpositive empirical label mean.
\end{lemma}
\begin{proof}
Set $A=\{M(z)>t\}$ and $p=\Pi(A)$. The low-output event has
\[
 \widehat R(z)\ge\widehat R(0)-2t\|y\|_S
               \ge\widehat R(0)-4at.
\]
The quadratic-gap inequality in Lemma~\ref{lem:canonical-prox} gives, on the high-output event,
\[
 \widehat R(z)+\lambda cM(z)
 \ge J_\lambda(u_\lambda)+\|z(S)-u_\lambda\|_S^2.
\]
Let $F$ be the posterior free energy. The entropy floor therefore implies
\begin{align*}
 F+\lambda e_0
 &\ge (1-p)\{\widehat R(0)-4at\}
       +p\{J_\lambda(u_\lambda)+\lambda B\}
       +\mathbb E_\Pi\mathbf1_A\|z(S)-u_\lambda\|_S^2.
\end{align*}
Extend $u_\lambda$ by zero at unobserved support points. Its sup norm is at most $2a$. The two variational comparisons give
\[
 F\le\min\left\{\widehat R(0)+s,
   J_\lambda(u_\lambda)+\lambda B+\zeta+
          \lambda[2a(4-c)+e_2]\right\}.
\]
For any two real numbers $A_0,A_1$ and nonnegative errors $r_0,r_1$,
\[
 \min(A_0+r_0,A_1+r_1)
 \le (1-p)A_0+pA_1+r_0+r_1.
\]
Subtracting the preceding inequalities bounds
$\mathbb E_\Pi\mathbf1_A\|z(S)-u_\lambda\|_S^2$ by the quantity under the square root. Meanwhile the low-output contribution to the mean has norm at most $t$. Cauchy--Schwarz controls the high-output contribution by that square root, proving the estimate.

Clipping contracts distance to $u_\lambda\in[-2a,2a]$ on the high event and cannot increase the low-event norm. Clipping the mean contracts distance to $p u_\lambda$. Finally the capped empirical-mean characterization from Lemma~\ref{lem:canonical-prox} proves the sign assertion.
\end{proof}

The baseline $B$ disappears from the approximation error. This cancellation is why a varying baseline cost does not itself protect a learner against the negative-singleton witness. It changes whether zero or a capped empirical predictor is favored, but both have the wrong sign on the same negative singleton observations when the true regression function is positive there.

\section{One regression problem for every mean-field-or-stronger scale}
\label{sec:all-scale-learning}
The baseline-plus-amplitude comparison removes the need to fix the mean-field variance prefactor. The following theorem uses one data distribution and one sample event for the entire scale interval.

The main learning statement is Theorem~\ref{thm:all-scale-learning}, stated at the start of the paper. We now prove its uniformity over the entire scale interval.

\begin{proof}[Proof of Theorem~\ref{thm:all-scale-learning}]
Write
\[
 \nu=Ns\in(0,1],\quad b=1-\nu,\quad
 \kappa=1-\nu+\nu\log\nu,\quad B=\kappa(H-1).
\]
The continuous value at $\nu=1$ is understood. The elementary inequality
\begin{equation}
 \kappa\ge b^2/2
 \label{eq:kappa-quadratic}
\end{equation}
follows by differentiating in $b$: $\frac{d}{db}\{b+(1-b)\log(1-b)\}=-\log(1-b)\ge b$ and both sides vanish at zero.

We use the joint entropy bound and uniform function comparison proved above. In a form sufficient here, for any fixed $t<2$ sufficiently close to two and any fixed threshold $u>0$, they give
\begin{align}
 2s\operatorname{KL}(\rho\|\mu)
 &\ge B\Pr_\rho(M(z)>u)+c_\nu\mathbb E_\rho M(z)-e_0,
 &c_\nu&=2\nu t,\label{eq:allscale-floor}\\
 e_0&\le C_t/H+C_t b(1+\log(1/u))
                  +C_t\log(2n)/N+4u,\label{eq:allscale-floor-error}
\end{align}
where $M(z)=\max_{x\in\mathcal U_D}|z(x)|$. Every $v:\mathcal U_D\to[-2a,2a]$ has a comparison law with arbitrary fixed RMS tolerance $\varepsilon>0$ and
\begin{align}
 2s\operatorname{KL}(\rho_v\|\mu)
 &\le B+4\|v\|_\infty+e_2,\label{eq:allscale-comparison}\\
 e_2&\le\kappa(d+1)+C_0a^2+
 C_0\left\{b(D+\log H)+\frac{(D+\log H)^2}{H}\right\}.
 \label{eq:allscale-comparison-error}
\end{align}
The constants are uniform over $0<\nu\le1$. A sufficient comparison width is $C_\varepsilon H^{C_\varepsilon}e^{C_\varepsilon D}$. All support-wide error bounds are simultaneous in the prescribed function. For the fixed choices below this threshold, and all additional polynomial width requirements, are absorbed by $N\ge n^C$ after increasing $C,H_0$.

The prior itself is centered with output variance at most $s\le1/N$, so its variational comparison is $\widehat R(0)+s$. Also $|Y|<2a$, and a comparison of RMS error $\varepsilon$ has empirical-loss error
\[
 \zeta\le8a\varepsilon+\varepsilon^2.
\]

Choose a sufficiently large absolute $\Lambda$, and divide temperatures at
\[
 \lambda_0=\frac{\Lambda a^2}{B+a}.
\]
This is only a proof division; the learner is allowed every positive temperature.

For $0<\lambda\le\lambda_0$, Lemma~\ref{lem:baseline-prox} puts each specified posterior prediction within
\begin{equation}
 u+\sqrt{4au+s+\zeta+
       \lambda[2a(4-c_\nu)+e_0+e_2]}
 \label{eq:allscale-low-error}
\end{equation}
of $p u_\lambda$ for some $p\in[0,1]$, where $u_\lambda$ is a capped empirical-mean vector. We show that the displayed distance can be made less than $a/1000$, uniformly in $\nu$.

The required uniform estimates are
\begin{equation}
 \frac1{B+a}\le\frac1a,\qquad
 \frac\kappa{B+a}\le\frac1{H-1},\qquad
 \frac b{B+a}\le\frac1{\sqrt{2a(H-1)}}.
 \label{eq:allscale-ratios}
\end{equation}
The last follows from~\eqref{eq:kappa-quadratic} by maximizing $x/[a+(H-1)x^2/2]$. Since
$4-c_\nu=4b+2\nu(2-t)$, these estimates imply that the contribution to the square in~\eqref{eq:allscale-low-error}, divided by $a^2$, has limsup as $H\to\infty$ at most
\begin{equation}
 C\Lambda\left\{(2-t)+a+\frac c{\sqrt a}+\frac{c^2}{a}+\frac ua\right\}
      +C\frac ua+C\frac\varepsilon a+C(\varepsilon/a)^2,
 \label{eq:allscale-limsup}
\end{equation}
with a fixed absolute $C$ (depending only on the fixed analytic constants). For clarity, the individual terms are controlled as follows: the first-layer cost contributes at most $\Lambda a^2(d+1)/(H-1)=o(a^2)$; the term $bD$ contributes at most $C\Lambda a^{3/2}c$; $D^2/H$ contributes at most $C\Lambda a c^2$; the coefficient discrepancy $4b$ and the terms $b\log H$ contribute $o(a^2)$; $C_0a^2$ contributes $C\Lambda a^3$. The terms $C_t/H$, $\log n/N$, and $s$ also vanish uniformly. All appearances of $C_t$ multiply terms vanishing with $H$, once $t$ is fixed.

First choose $\Lambda$ as large as needed for the high-temperature estimates below. Then choose $2-t>0$ small, $a>0$ small within the analytic amplitude range, $c>0$ sufficiently small compared with $\sqrt a/\Lambda$, and $u,\varepsilon>0$ sufficiently small compared with $a/\Lambda$. The right side of~\eqref{eq:allscale-limsup} can thereby be made less than any prescribed positive constant, in particular less than $10^{-8}$. Increase $H_0,C$ to make~\eqref{eq:allscale-low-error} less than $a/1000$ at every stated width.

For $\lambda>\lambda_0$, the prior comparison gives
\[
 F_S\le\widehat R(0)+s\le5a^2,
 \qquad \mathbb E_\Pi\widehat R(z)\le F_S,
 \qquad \mathbb E_\Pi\|z(S)\|_S^2\le18a^2,
\]
after $1/N\le a^2$. If $B\ge a$, drop the nonnegative amplitude term in~\eqref{eq:allscale-floor}. With $p=\Pi(M(z)>u)$,
\[
 p\le\frac{5a^2}{\lambda B}+\frac{e_0}B
       \le\frac{10}{\Lambda}+\frac{e_0}B.
\]
Uniformly in this range, $e_0/B\le4u/a+o(1)$: the $b$ term is bounded using $B\ge a$ and~\eqref{eq:kappa-quadratic}, giving $b/B\le C/\sqrt{aH}$, and the remaining terms are immediate. Therefore
\[
 \|\mathbb E_\Pi z(S)\|_S
 \le u+\sqrt{18a^2p}<a/1000
\]
by the preceding choices of $\Lambda,u$ and sufficiently large $H_0$. The same inequality holds for both clipped prediction conventions.

If instead $B<a$, then $b\le\sqrt{2a/(H-1)}$, so $c_\nu\ge3$ for large $H$. Dropping the baseline in~\eqref{eq:allscale-floor} yields
\[
 \mathbb E_\Pi M(z)
 \le\frac{5a^2}{c_\nu\lambda}+\frac{e_0}{c_\nu}
 \le\frac{10a}{3\Lambda}+\frac{4u+o(1)}3<a/1000.
\]
This also controls all stated prediction conventions. Thus every high-temperature prediction is close to the zero vector, which is a capped empirical-mean reference.

On the sample event used in Theorem~\ref{thm:canonical-learning}, at least $n/10$ observations are singleton inputs with negative noise sign. This event has probability at least $9/10$ for all sufficiently large $H$. At each such input, the empirical mean is negative while $g(X)>a/4$. Every reference vector just obtained is nonpositive there. Its empirical error is at most $a/1000$, so the common witness $w_i=-1$ on these indices and zero elsewhere satisfies
\[
 \langle w,\widehat f-g\rangle_S
 \ge a/40-a/1000>a/64.
\]
All reference vectors lie in $[-2a,2a]$ and have the same sign property. Their convex combinations retain these properties; averaging the error vectors cannot increase their norm, and clipping after averaging contracts distance to the bounded reference. The witness therefore applies to arbitrary well-defined mixtures over every allowed width, scale, and temperature.

Because singleton atoms have population probability $1/n$, the corresponding population witness has squared norm at most one and the same pairing. Cauchy--Schwarz gives excess squared loss at least $a^2/4096$. The bounded-width teacher construction is the same as in Theorem~\ref{thm:canonical-learning}. This proves the theorem.
\end{proof}

The result is an obstruction to a sample bound polynomial jointly in input dimension, learner depth, and the complexity of an approximating teacher for this unrestricted Gaussian learner. It is consistent with the fixed-depth approximation theorems. The next theorem gives a uniform quantitative refinement below criticality, interpolating between $\exp(\Omega(\sqrt H))$ and $\exp(\Omega(H))$ as the variance gap grows.

\section{The transition from square-root to linear depth}
\label{sec:prefactor-window}
The common obstruction can be sharpened when the variance prefactor is
known to stay below one by a possibly depth-dependent gap. The transition
occurs at a gap of order $H^{-1/2}$; no fixed separation from one is needed.

\begin{theorem}[A uniform critical-window obstruction]
\label{thm:prefactor-window}
There are absolute constants $a,c,C>0$ and $H_0$ with the following
property. For any integer $H\ge H_0$ and any prescribed
$0\le\Delta<1$, put
\[
 D=\left\lfloor c(\sqrt H+\Delta H)\right\rfloor,
 \qquad d=D+2,\qquad n=4\cdot2^D.
\]
Let $X$ be uniform on $\mathcal U_D$, and let
\[
 Y=\frac a4(2+\tanh X_1)+a\epsilon,
 \qquad \Pr(\epsilon=1)=\Pr(\epsilon=-1)=\tfrac12,
\]
with independent noise. With probability at least $9/10$ over $n$ IID
observations, every Gaussian learner of Theorem~\ref{thm:all-scale-learning}
with
\[
 N\ge n^C,\qquad 0<Ns\le1-\Delta,\qquad\lambda>0
\]
has excess population squared loss at least $a^2/4096$.
The same statement covers its raw posterior mean, the two clipped
prediction conventions, and arbitrary well-defined data-dependent
mixtures over the allowed widths, scales and temperatures.
All constants are independent of $\Delta$. The conditional mean has
width two, one hidden layer, and uniformly bounded coefficients.
\end{theorem}

Thus the sample obstruction is
$\exp(\Omega(\sqrt H+\Delta H))$. At $\Delta=0$ it is the all-scale
square-root-depth obstruction. At a fixed positive gap it is exponential
in depth. In the intermediate range it gives a quantitative crossover,
while input dimension grows as $\sqrt H+\Delta H$. It is not a
fixed-input-dimension lower bound.

\begin{proof}
Use the notation of the proof of Theorem~\ref{thm:all-scale-learning}:
$\nu=Ns$, $b=1-\nu\ge\Delta$,
$\kappa=1-\nu+\nu\log\nu$, and $B=\kappa(H-1)$.
The elementary bound $\kappa\ge b^2/2$ gives, for $H\ge2$,
\begin{align}
 \frac{bD}{B+a}
 &\le c\left(\frac1{\sqrt a}+4\right),
 \label{eq:window-linear-ratio}\\
 \frac{D^2}{H(B+a)}
 &\le 2c^2\left(\frac1a+4\right),
 \label{eq:window-quadratic-ratio}\\
 \frac{\kappa(d+1)}{B+a}
 &\le 2c(H^{-1/2}+\Delta)+\frac6H.
 \label{eq:window-first-ratio}
\end{align}
To verify the first inequality, split $D\le c\sqrt H+c\Delta H$.
Maximizing $b/[a+(H-1)b^2/2]$ bounds the first contribution by
$c\sqrt{H/[2a(H-1)]}\le c/\sqrt a$.
If $\Delta>0$, the second is at most
$2c(\Delta/b)H/(H-1)\le4c$; if $\Delta=0$ it vanishes.
For the second inequality use
$D^2\le2c^2(H+\Delta^2H^2)$ and
$B\ge\Delta^2(H-1)/2$; again the term containing $\Delta$ is zero
when $\Delta=0$.
The third follows from $\kappa/(B+a)\le1/(H-1)$ and $d+1=D+3$.

The uniform function comparison applies because
$D+\log H\le c_0H$ after choosing $c$ small enough and then $H_0$
large enough. Its width requirement
$C_\varepsilon H^{C_\varepsilon}e^{C_\varepsilon D}$ is absorbed by
$N\ge n^C$ uniformly in $\Delta$: $D\ge c\sqrt H/2$ for large $H$.
The same observation absorbs every finite-width concentration error.

Retain the same temperature division
$\lambda_0=\Lambda a^2/(B+a)$. The low-temperature proximity bound is
exactly \eqref{eq:allscale-low-error}. In its square, divide by $a^2$.
The only previously used estimates that involved $D$ were the
first-layer term $\kappa(d+1)$, the growing-prefix term $bD$, and
the quadratic term $D^2/H$. Equations
\eqref{eq:window-linear-ratio}--\eqref{eq:window-first-ratio}
replace those three estimates. The $\log H$ terms are still uniformly
vanishing, using
\[
 \frac b{B+a}\le\frac1{\sqrt{2a(H-1)}},\qquad
 \frac1{B+a}\le\frac1a.
\]
The resulting uniform limiting error is at most
\[
 C\Lambda\left\{(2-t)+a+c+\frac c{\sqrt a}
       +\frac{c^2}a+\frac ua\right\}
   +C\frac ua+C\frac\varepsilon a+C(\varepsilon/a)^2.
\]
Choose $\Lambda$ first, then $t<2$ sufficiently close to two,
then $a>0$ small, then $c>0$ small, and finally $u,\varepsilon>0$
small. These choices make this bound arbitrarily small, independently
of $\Delta$. Increasing $H_0,C$ gives the same proximity $a/1000$
as before.

The high-temperature argument does not use $D$ except through
$\log n/N$, which still tends uniformly to zero at the stated width.
The alternatives $B\ge a$ and $B<a$ therefore give the same proximity
to zero, with the same constants. Finally the singleton-input noise
witness has the same probability and population conversion on the
$n$-point uniform law for every $D$. It gives excess $a^2/4096$ for
all allowed predictions, and is preserved by the stated mixtures
and clipping conventions exactly as in the all-scale proof.
\end{proof}

\begin{corollary}[The same crossover for the observation-model extensions]
\label{cor:prefactor-window-extensions}
The order $n=\exp(\Omega(\sqrt H+\Delta H))$, the width threshold $N\ge n^C$, and uniformity over $0<Ns\le1-\Delta$ also hold for each of the following extensions, with their existing prediction conventions and adjusted absolute constants:
\begin{enumerate}
\item Final clipping or tanh inside both likelihood and prediction, with bounded independent noise, as in Corollary~\ref{cor:all-scale-bounded-likelihood}.
\item Arbitrary independent Gaussian affine output components under the raw likelihood, with bounded independent noise, as in Theorem~\ref{thm:all-scale-affine-likelihood}.
\item Correctly specified Gaussian observation noise and arbitrary independent Gaussian affine output components under the raw likelihood, as in Corollary~\ref{cor:all-scale-gaussian-affine}.
\end{enumerate}
Constants are uniform in $\Delta$. Mixtures are permitted exactly within the prediction families stated in those results. This does not add clipping of an arbitrary affine prediction or a bounded link after an affine skip.
\end{corollary}
\begin{proof}
In the bounded-output comparison, inverse tanh enlarges the fixed amplitude cap only from $2a$ to $3a$. In the Gaussian-noise comparison, clipping the empirical atom means enlarges it to $32a$. Thus their comparison errors have the same form as \eqref{eq:allscale-comparison-error}, with larger fixed coefficients on $a^2$, $b(D+\log H)$ and $(D+\log H)^2/H$. The three ratios \eqref{eq:window-linear-ratio}--\eqref{eq:window-first-ratio} therefore control every dimension-dependent term in their low-temperature arguments as well. The fixed tolerance and amplitude constants can be chosen smaller in the same order as above, independently of $\Delta$.

For affine output components, Gaussian integration replaces the sample norm by a contraction norm. Its proximity estimates are uniform in that contraction and depend on $D$ only through the same comparison error. The high-temperature bounds and Gaussian-noise clipping error are unchanged. The remaining input Gram and affine sign-pattern failure bounds have the forms
\[
 C(D+3)^2 e^{-c n/(D+3)^4},\qquad
 \exp\{C(D+5)\log(2n)-c n\},
\]
and tend uniformly to zero because $n=4\cdot2^D$ and $D\ge c\sqrt H/2$. Occupancy, bounded-noise alignment and Gaussian clipping events also tend uniformly to probability one. Their common witnesses and empirical-to-population conversions are unchanged. The finite-width comparison requirement remains $C_\varepsilon H^{C_\varepsilon}e^{C_\varepsilon D}$ and is absorbed by $N\ge n^C$. This proves each extension with exactly its previously specified prediction scope.
\end{proof}

\section{Gaussian label noise throughout the scale range}
\label{sec:all-scale-Gaussian-noise}
The scale-uniform conclusion also holds when the labels have exactly a Gaussian conditional distribution.

\begin{corollary}[One Gaussian-noise problem for every allowed prior scale]
\label{cor:all-scale-Gaussian-noise}
For suitable absolute constants $a,c,C,H_0>0$, the setting of Theorem~\ref{thm:all-scale-learning}, with labels instead given by
\[
 Y=g(X)+aZ,\qquad Z\sim\mathcal N(0,1)\text{ independently of }X,
 \qquad g(X)=\frac a4(2+\tanh X_1),
\]
has, with probability at least $9/10$,
\[
 R(\widehat f)\ge a^2+a^2/512^2
\]
simultaneously for every $N\ge n^C$, $0<s\le1/N$, every positive Gaussian-likelihood variance, and every prediction or mixture listed in Theorem~\ref{thm:all-scale-learning}. In particular the true observation variance $a^2$ is included.
\end{corollary}
\begin{proof}
We spell out the changes to the scale-uniform proof. Write $B,c_\nu,e_0,e_2$ as in~\eqref{eq:allscale-floor}--\eqref{eq:allscale-comparison-error}, but use comparison cap $A=16a$. Enlarging the fixed coefficient in $C_0a^2$ preserves the same error form. Put $r_0=\widehat R_S(0)$, and work on the sample events
\begin{equation}
 r_0\le6a^2,\qquad
 T_S:=\frac1n\sum_i(|b_i|-A)_+^2\le2^{-50}a^2,
 \label{eq:allscale-Gaussian-events}
\end{equation}
where $b_i$ are empirical means within repeated-input groups. The cap-error argument in Section~\ref{sec:canonical-Gaussian-noise} gives
\[
 J_\lambda(v_\lambda)-J_\lambda(u_\lambda)\le T_S,
 \qquad v_\lambda=\operatorname{clip}(u_\lambda,[-A,A]).
\]
Here $u_\lambda$ minimizes $J_\lambda(u)=\widehat R_S(u)+\lambda c_\nu\|u\|_\infty$ on repeated-input sample vectors. Although $u_\lambda$ need not be bounded by $A$, $v_\lambda$ is an admissible comparison. Its approximation incurs empirical-risk error at most
$\zeta\le38a\varepsilon+\varepsilon^2$ on~\eqref{eq:allscale-Gaussian-events}.

The proof of Lemma~\ref{lem:baseline-prox} now gives, for $p=\Pi\{M(z)>u\}$,
\begin{align}
 \|\E_\Pi z(S)-p u_\lambda\|_S
 \le u+\bigl[2\sqrt6au+s+T_S+\zeta
       +\lambda\{A(4-c_\nu)+e_0+e_2\}\bigr]^{1/2}.
 \label{eq:allscale-Gaussian-low}
\end{align}
Indeed the low-output cross term is bounded by $2u\|Y\|_S\le2\sqrt6au$. On the high-output event the same strong-convexity gap applies. In the upper comparison, clipping $u_\lambda$ to $v_\lambda$ costs $T_S$, and
\[
 4\|v_\lambda\|_\infty-c_\nu\|u_\lambda\|_\infty
 \le(4-c_\nu)A.
\]
These are exactly the two changes to the baseline-cancellation argument.

Divide temperatures again at $\lambda_0=\Lambda a^2/(B+a)$. The three uniform ratios in~\eqref{eq:allscale-ratios} still control every term in~\eqref{eq:allscale-Gaussian-low}. Its squared error divided by $a^2$ has the same limsup bound~\eqref{eq:allscale-limsup}, with larger absolute constants and the extra term $2^{-50}$. Choose $\Lambda$ large, then $2-t$, $a$, $c$, $u/a$, and $\varepsilon/a$ sufficiently small, in that order. For large enough $H$ and $N\ge n^C$, every $\lambda\le\lambda_0$ then gives empirical error less than $a/10000$. Clipped predictions satisfy the same error bound relative to the corresponding clipped capped references; these references remain nonpositive wherever $b_i<0$.

For $\lambda>\lambda_0$, the prior comparison gives $F_S\le7a^2$ and
$\E_\Pi\|z(S)\|_S^2\le26a^2$. If $B\ge a$, the threshold-event probability obeys
\[
 p\le14/\Lambda+e_0/B\le14/\Lambda+4u/a+o(1),
\]
uniformly in the scale. Thus the raw mean has empirical norm at most
$u+\sqrt{26a^2p}<a/10000$; the same holds for both clipped conventions. If $B<a$, then $c_\nu\ge3$ for large $H$, and the amplitude floor yields
\[
 \E_\Pi M(z)\le14a/(3\Lambda)+(4u+o(1))/3<a/10000.
\]
The choices can satisfy the low- and high-temperature requirements simultaneously.

Finally, the Gaussian sample estimates in Section~\ref{sec:canonical-Gaussian-noise} establish~\eqref{eq:allscale-Gaussian-events} together with at least $n/100$ singleton observations having $Z\le-1$, with joint probability at least $9/10$ at large $H$. At these inputs the empirical labels are negative and $g>a/4$. Every preceding reference is nonpositive there, so the same witness gives
\[
 \langle w,\widehat f-g\rangle_S\ge a/400-a/10000>a/512.
\]
References are bounded on the finite observed sample by its largest label magnitude. Convexity, clipping contraction, and the common sign property therefore preserve the witness for all well-defined mixtures. Its population translation proves squared excess risk at least $a^2/512^2$.
\end{proof}

The correctly specified likelihood variance corresponds to $\lambda=a^2/(ns)$ in the displayed posterior of Theorem~\ref{thm:all-scale-learning}. Its inclusion is automatic because the sample event covers every positive $\lambda$ and every allowed $s$ simultaneously.

\section{Bounded outputs inside the likelihood at every scale}
\label{sec:all-scale-bounded-likelihood}
The common obstruction is not specific to an unbounded linear output in the observation model. Both final clipping and a final tanh preserve it.

\begin{corollary}[Final clipping or tanh in the observation model]
\label{cor:all-scale-bounded-likelihood}
In Theorem~\ref{thm:all-scale-learning}, let $\psi$ be either clipping to $[-1,1]$ or tanh. Form the posterior using the actual output $\psi(z_w)$:
\[
 \Pi^\psi(dw)\propto
       \exp\{-\widehat R_S(\psi(z_w))/(2\lambda s)\}\,\mu_{N,s}(dw),
 \qquad \widehat f_\psi(x)=\E_{\Pi^\psi}\psi(z_w(x)).
\]
After adjusting the absolute constants $a,c,C,H_0$, the same input and label law as in that theorem has, with probability at least $9/10$, the simultaneous lower bound
\[
 R(\widehat f_\psi)\ge a^2+a^2/4096
\]
for all $N\ge n^C$, all $0<s\le1/N$, and all $\lambda>0$. The event can be chosen common to both bounded-output likelihoods and the original linear-output likelihood. It therefore also covers arbitrary data-dependent mixtures of these posterior predictions and the prediction conventions in Theorem~\ref{thm:all-scale-learning}, including clipping after mixing.

The conditional mean remains $g(x)=a(2+\tanh X_1)/4$. With a final clipping output it has the same width-two shallow implementation. With a final tanh output it has a width-four implementation with a fixed coefficient bound on the support, at every hidden depth. Thus the bounded observation model still contains a fixed-complexity teacher realizing the true conditional mean.
\end{corollary}
\begin{proof}
Both output maps are odd, vanish at zero, and are one-Lipschitz. Put
\[
 M_\psi(w)=\max_{x\in\mathcal U_D}|\psi(z_w(x))|.
\]
Then $M_\psi\le M$. Since the depth baseline $B$ and amplitude coefficient $c_\nu$ are nonnegative, the joint entropy floor implies
\[
 2s\operatorname{KL}(\rho\|\mu)
 \ge B\Pr_\rho(M_\psi>u)+c_\nu\E_\rho M_\psi-e_0.
\]
The prior prediction remains centered and has pointwise variance at most $s$, so the same zero comparison applies to its actual squared loss.

For every comparison function $v:\mathcal U_D\to[-2a,2a]$, use raw target $v$ for clipping and raw target $\operatorname{arctanh}v$ for a final tanh. For sufficiently small fixed $a$,
\[
 \|\operatorname{arctanh}v\|_\infty
       \le\|v\|_\infty+Ca^3\le3a.
\]
The uniform function comparison therefore has normalized entropy at most
\[
 B+4\|v\|_\infty+e_2',\qquad
 e_2'\le\kappa(d+1)+C'a^2+
 C'\left\{(1-\nu)(D+\log H)+\frac{(D+\log H)^2}{H}\right\}.
\]
Contraction by $\psi$ preserves its prescribed RMS supremum accuracy. In particular, its expected actual empirical loss is at most $\widehat R_S(v)+8a\varepsilon+\varepsilon^2$, exactly as in the linear-output proof.

The baseline-plus-amplitude lemma applies to the actual function $\psi(z)$ with these comparisons and entropy floor. The single temperature cutoff and all uniform parameter estimates in the proof of Theorem~\ref{thm:all-scale-learning} are unchanged after increasing the fixed constants in $e_2$. The high-temperature branch uses only the zero comparison, the entropy floor, and a second-moment bound for the actual prediction, so it is valid here as well. Thus every stated prediction is close to zero or to a nonnegative multiple of a capped empirical-mean vector, with the same error tolerance. On every negative singleton observation those reference predictions are nonpositive, whereas $g>a/4$. The same common witness proves the claimed excess and all mixture assertions. Since the reference vectors lie in $[-2a,2a]$, clipping a mixture preserves their approximation.

It remains to verify realizability for the final-tanh output while retaining the same conditional mean. Let $T_j$ run through $(-3/4,-1/4,1/4,3/4)$, and form the four-feature matrix whose $j$th column is
\[
 F_j=\bigl(\tanh1,\ \tanh(20(T_j+1/2)),\
                 \tanh(20T_j),\ \tanh(20(T_j-1/2))\bigr)^{\mathsf T}.
\]
This matrix is invertible. Replacing the last three rows by their sign limits gives the usual constant-and-three-step basis on four ordered points, whose inverse has operator norm less than two. The actual matrix differs from it in Frobenius norm by less than $10^{-3}$ because each nonzero threshold distance is at least $1/4$ and $1-\tanh5<10^{-4}$. Hence invertibility follows from the Neumann criterion.

Write $b_j=\operatorname{arctanh}(g(T_j))$ and choose $c^{\mathsf T}=b^{\mathsf T}F^{-1}$. The hidden features $F_j$ followed by readout $c$ and final tanh give $g(T_j)$ exactly. To preserve these features at arbitrary hidden depth, use the fixed matrix
\[
 W=\operatorname{arctanh}(F)\,F^{-1},
\]
where inverse tanh acts entrywise. Then $\tanh(WF)=F$. All first-layer parameters, entries of $W$, and readout coefficients are bounded by fixed constants independent of $H,D,N$; the readout coefficients are $O(a)$. No internal biases are needed for this relay. This proves the width-four teacher assertion.
\end{proof}

For the bounded likelihood, the prediction here is the mean of the bounded network output. A final tanh applied to the mean of the raw readout is a different convention and is not asserted by this corollary. The affine-output extension below concerns a raw Gaussian likelihood; it is not asserted for these bounded-output likelihoods.

\section{Arbitrary Gaussian affine output components at every scale}
\label{sec:all-scale-affine-likelihood}
An independently regularized output bias or direct linear input map does not remove the common obstruction. The proof uses the geometry of a capped affine fit and a uniform sign-pattern bound on the singleton observations.

\begin{theorem}[The all-scale obstruction with affine output skips]
\label{thm:all-scale-affine-likelihood}
The conclusion of Theorem~\ref{thm:all-scale-learning} holds, after adjusting its absolute constants and replacing its excess lower bound by $2^{-60}a^2$, when the learner predicts
\[
 z_w(x)+b^{\mathsf T}(1,x),
\]
where $b$ has any independent centered Gaussian prior with a finite positive semidefinite covariance matrix. The statement concerns ordinary raw posterior means under the Gaussian likelihood. It is simultaneous over every such covariance, every $N\ge n^C$, every $0<s\le1/N$, and every positive likelihood temperature, and it survives arbitrary well-defined data-dependent mixtures of these predictions. Singular covariance, including no affine component, is permitted.

The data law and conditional mean are the same as in Theorem~\ref{thm:all-scale-learning}; in particular the conditional mean is a fixed width-two shallow tanh function. The input set lies on the unit sphere and the sample count remains $\exp(\Omega(\sqrt H))$.
\end{theorem}

\begin{lemma}[The geometry of the affine proximal reference]
\label{lem:affine-prox-geometry}
Let $b$ be the empirical atom-mean label vector and let
\[
 B=(I+\Phi\Sigma\Phi^{\mathsf T}/(n\lambda s))^{-1},
\]
where $\Phi$ is the affine training-feature matrix and $\Sigma\succeq0$ is finite. For $c>0$, let $u$ minimize
\[
 J_B(u)=\|B^{1/2}(b-u)\|_n^2+\lambda c\|u\|_\infty
\]
over sample vectors constant on repeated inputs. Put
\[
 f_1=(I-B)b+Bu,\qquad f_0=(I-B)b.
\]
Then $\|u\|_\infty\le\|b\|_\infty$. There are an affine function $h$ and a cap $t\ge0$ for which
\[
 f_1=h+\operatorname{clip}(b-h,[-t,t]).
\]
Both reference vectors obey the ordinary sector inequality
\[
 \langle b,f_j\rangle_n\ge\|f_j\|_n^2,
 \qquad j\in\{0,1\}.
\]
\end{lemma}
\begin{proof}
The candidate $u=b$ has objective $\lambda c\|b\|_\infty$. Every candidate with larger sup norm has strictly greater penalty alone, so it cannot minimize. Gaussian least-squares completion gives
\[
 \|B^{1/2}(b-u)\|_n^2
 =\min_v\{\|b-u-\Phi\Sigma^{1/2}v\|_n^2
                                      +\lambda s\|v\|^2\}.
\]
At the optimum set $h=\Phi\Sigma^{1/2}v=(I-B)(b-u)$ and $t=\|u\|_\infty$. With $h,t$ fixed, the least-squares minimizing $u$ is exactly the coordinatewise clipping of $b-h$ to $[-t,t]$. Repeated observations have identical $b,h$, so this minimizer belongs to the required atom-constant subspace. This proves the asserted form of $f_1=h+u$.

Varying the optimizer along scalar multiples of $u$ gives
\[
 \langle u,B(b-u)\rangle_n
                   =\tfrac12\lambda c\|u\|_\infty\ge0.
\]
Writing $v=b-u$, direct algebra yields
\[
 \langle b,f_1\rangle_n-\|f_1\|_n^2
 =\langle u,Bv\rangle_n+\langle v,(B-B^2)v\rangle_n\ge0,
\]
because $0\prec B\preceq I$. The same inequality for $f_0$ follows from $0\preceq I-B\preceq I$.
\end{proof}

\begin{lemma}[Every capped affine fit has a fixed empirical error]
\label{lem:capped-affine-error}
For the sample law of Theorem~\ref{thm:all-scale-learning}, with probability tending to one as $H\to\infty$, every affine function $h$ and every $t\ge0$ satisfy
\[
 \left\|h+\operatorname{clip}(b-h,[-t,t])-g(S)\right\|_n^2
                          \ge2^{-24}a^2.
\]
The event depends only on the observed sample, not on a prior or a temperature.
\end{lemma}
\begin{proof}
The population squared distance of $\tanh T$ from affine functions of the inputs is
\[
 r_*=
 \frac{\{(3/4)\tanh(1/4)-(1/4)\tanh(3/4)\}^2}
      {2\{(1/4)^2+(3/4)^2\}}>1/4096.
\]
Indeed the nuisance coordinates are independent centered signs, the last sphere coordinate is even in $T$, and the only odd affine direction is $T$. The constant part of $g=a/2+(a/4)\tanh T$ is affine, so its squared affine residual is $a^2r_*/16>a^2/65536$.

Use the whitened affine features $(1,T,z,\sqrt{3/4-T^2})$, which span the same affine functions as the actual sphere coordinates. Their population Gram matrix has a fixed positive minimum eigenvalue: the nuisance block is the identity, the $T$ variance is $5/16$, and the constant/even-coordinate block is a fixed invertible two-by-two matrix. Bounded-entry concentration for this Gram matrix and for the moments of $g/a$ therefore gives
\[
 \inf_{h\text{ affine}}\|h(S)-g(S)\|_n\ge a/512
\]
with probability tending to one. For an explicit sufficient estimate, entry errors $O((D+3)^{-2})$ preserve the Gram inverse and the positive Schur complement; a union bound has failure at most $C(D+3)^2e^{-c n/(D+3)^4}$. Work on this event and put $\delta=a/1024$.

If $t\le\delta$, clipping changes $h(S)$ in empirical norm by at most $\delta$. Hence the prediction's distance from $g(S)$ is at least $a/1024$, which is stronger than the claimed lower bound.

Suppose now $t\ge\delta$. With probability tending to one, at least $n/3$ inputs occur exactly once. Conditional on all input observations, their independent label-noise signs remain fair. For a singleton at $x$, write $d_x=h(x)-g(x)$. If $d_x\ge0$ and the noise sign is positive, then
\[
 h(x)+\operatorname{clip}(g(x)+a-h(x),[-t,t])-g(x)
   =\operatorname{clip}(a,[d_x-t,d_x+t])\ge\min(a,t)\ge\delta.
\]
If $d_x<0$ and the noise sign is negative, the analogous error is at most $-\delta$.

The patterns $\operatorname{sgn}(h(x)-g(x))$ on the singleton inputs, with zero assigned the positive sign, range over at most $(2n)^{D+5}$ possibilities. This follows, for example, by viewing them as halfspace labels of the augmented vectors $((1,x),g(x))$ and applying the elementary hyperplane-arrangement bound. For each fixed pattern, a Chernoff bound shows that at least one third of the singleton noise signs match it except on probability $e^{-c n}$. A union bound therefore gives this matching property simultaneously for every affine $h$, with failure at most $\exp\{C(D+5)\log(2n)-c n\}$, which tends to zero. At least $n/9$ singleton errors consequently have magnitude at least $\delta$, uniformly for all $h$ and all $t\ge\delta$. The squared empirical error is at least $\delta^2/9>2^{-24}a^2$.
\end{proof}

\begin{proof}[Proof of Theorem~\ref{thm:all-scale-affine-likelihood}]
Integrating the affine Gaussian block gives the exact marginal neural posterior with loss
\[
 \widehat R_B(z)=\|B^{1/2}(y-z(S))\|_n^2,
\]
where $B$ is as above. Its determinant factor is independent of the neural weights and cancels from the marginal posterior. If $m_z$ is the marginal neural mean, the total posterior mean on the sample is
\[
 f=(I-B)y+Bm_z=(I-B)b+Bm_z.
\]
The second identity holds because the within-atom residual $y-b$ is orthogonal to every affine feature and hence is fixed by $B$.

The baseline-plus-amplitude argument applies in the $B$ inner product. The prior zero comparison is at most $\widehat R_B(0)+s\le4a^2+s$. The joint entropy floor is unchanged, since it concerns the neural prior. If the raw support maximum is at most the threshold, its $B$-norm sample vector is also at most that threshold. Lemma~\ref{lem:affine-prox-geometry} gives a proximal reference $u$ with sup norm at most $2a$; extending it by zero off the observed atoms permits exactly the same uniform finite-function comparison. Contraction by $B^{1/2}$ does not increase its RMS error. The quadratic-gap estimate holds in the $B$ norm.

It follows, using the same single temperature cutoff and uniform estimates as in Theorem~\ref{thm:all-scale-learning}, that the constants can be chosen so that for each low-temperature learner and some $p\in[0,1]$,
\[
 \|B^{1/2}(m_z-pu)\|_n\le2^{-28}a.
\]
The requested smaller fixed tolerance is obtained by the same choices of the fixed amplitude, support coefficient, width threshold, and likelihood cutoff constants. All estimates are uniform in $B$, since $0\prec B\preceq I$. For the high-temperature branch, its indicator argument uses $\E\|B^{1/2}z(S)\|_n^2\le18a^2$, and its amplitude argument directly bounds the neural support maximum. These give $\|Bm_z\|_n\le2^{-28}a$, uniformly in $B$. Therefore every full posterior mean is within $2^{-28}a$ of either $f_0$, or the convex combination $(1-p)f_0+pf_1$, from Lemma~\ref{lem:affine-prox-geometry}.

Let $w=b-2g(S)$. On an event of probability tending to one, the noise alignment satisfies
\[
 |\langle b-g(S),g(S)\rangle_n|\le2^{-27}a^2
\]
by conditional Hoeffding. Both references have squared empirical error at least $2^{-24}a^2$ by Lemma~\ref{lem:capped-affine-error}. Their sector inequalities imply
\begin{align*}
 \langle w,f_j-g(S)\rangle_n
 &\ge\|f_j-g(S)\|_n^2
                -\langle b-g(S),g(S)\rangle_n\\
 &\ge2^{-25}a^2.
\end{align*}
This same linear witness holds for every convex combination of the references. Since $\|w\|_n\le2a$, the posterior approximation error reduces it by at most $2^{-27}a^2$. Thus every allowed posterior mean satisfies
\[
 \langle w,f-g(S)\rangle_n\ge2^{-26}a^2.
\]
The witness is independent of the affine covariance, width, scale, and temperature, and therefore survives arbitrary mixtures.

For the population conversion, define
\[
 h_S(x)=\sum_{i:X_i=x}\{Y_i-2g(x)\}.
\]
The uniform input law assigns each atom mass $1/n$, so $\langle h_S,f-g\rangle_{L^2(P_X)}=\langle w,f-g(S)\rangle_n$. On the occupancy event $n^{-1}\sum_xk_x^2\le4$, and $|Y_i-2g(X_i)|\le2a$, we have $\|h_S\|_{L^2(P_X)}^2\le16a^2$. Cauchy--Schwarz gives excess risk at least $2^{-56}a^2$, leaving slack for the stated $2^{-60}a^2$. All sample events have total failure below one tenth after increasing the fixed depth threshold, because $n=4\cdot2^{\lfloor c\sqrt H\rfloor}$ dominates every polynomial in $D$. This proves the theorem.
\end{proof}

The affine theorem concerns raw Gaussian observation models and raw posterior means. Clipping a general affine prediction can create a nonlinear component, and this proof does not automatically extend to that convention or to a bounded link applied after the affine skip.

\section{Correctly specified Gaussian noise with affine output skips}
\label{sec:all-scale-gaussian-affine}
The all-scale affine result also holds when the observation-noise family is correctly specified. Only a proof comparison is clipped; the learner uses the original Gaussian labels.

\begin{corollary}[Gaussian-noise obstruction with arbitrary affine covariance]
\label{cor:all-scale-gaussian-affine}
Replace the labels in Theorem~\ref{thm:all-scale-affine-likelihood} by
\[
 Y=g(X)+aG,\qquad G\sim N(0,1)\text{ independent of }X,
 \qquad g(X)=\frac a4(2+\tanh X_1).
\]
After adjusting its absolute constants and replacing its excess lower bound by $2^{-64}a^2$, all its conclusions hold for ordinary raw posterior means and their arbitrary data-dependent mixtures. In particular the same sample event covers every width $N\ge n^C$, every $0<s\le1/N$, every positive likelihood temperature, and every independent centered Gaussian affine-output covariance. The likelihood family contains the true Gaussian observation variance $a^2$.
\end{corollary}

\begin{lemma}[A bounded comparison for an unbounded label vector]
\label{lem:clipped-affine-prox-comparison}
Let $B$ be any positive definite contraction preserving the atom-constant subspace, let $b$ belong to that subspace, and put $b_A=\operatorname{clip}(b,[-A,A])$. For $c,\lambda>0$ define
\[
 J_b(u)=\|B^{1/2}(b-u)\|_n^2+\lambda c\|u\|_\infty,
 \qquad u_b=\arg\min J_b,
\]
and define $u_A$ using $b_A$ in place of $b$. Then
\[
 \|u_A\|_\infty\le A,
 \qquad
 J_b(u_A)-J_b(u_b)\le\|B^{1/2}(b-b_A)\|_n^2.
\]
\end{lemma}
\begin{proof}
The cap bound follows by comparing with $u=b_A$. Strong convexity of $J_{b_A}$ gives
\[
 J_{b_A}(u_A)-J_{b_A}(u_b)
                \le-\|B^{1/2}(u_A-u_b)\|_n^2.
\]
Expanding the two squared losses after replacing $b_A$ by $b$ changes their difference by
$-2\langle B^{1/2}(b-b_A),B^{1/2}(u_A-u_b)\rangle_n$. If $r=\|B^{1/2}(b-b_A)\|_n$ and $d=\|B^{1/2}(u_A-u_b)\|_n$, the resulting upper bound is $-d^2+2rd\le r^2$.
\end{proof}

\begin{proof}[Proof of Corollary~\ref{cor:all-scale-gaussian-affine}]
Use the same input law and let $b$ denote the empirical atom means. All constants below are fixed independently of depth and prior parameters.

\emph{Sample bounds and clipping error.}
With probability tending to one, the empirical squared Gaussian-noise average is at most $2a^2$. Since $|g|<3a/4$, this gives
\[
 \|y\|_n\le3a,\qquad \|b\|_n\le3a,
 \qquad \|b-2g(S)\|_n\le3a.
\]
The occupancy events from the previous proof continue to give at least $n/3$ singleton inputs and $n^{-1}\sum_xk_x^2\le4$.
Set $A=32a$ and $b_A=\operatorname{clip}(b,[-A,A])$. Conditional on the inputs, an atom of count $k\ge1$ has mean $b_x=g(x)+aZ_x/\sqrt k$ with independent standard Gaussians $Z_x$. Therefore
\[
 k(b_x-b_{A,x})^2\le a^2(|Z_x|-31)_+^2.
\]
For $W=(|Z|-31)_+^2$, Gaussian tails give
$\Pr(W>u)\le2e^{-31^2/2}e^{-u/2}$ and $\E e^{W/4}\le1+2e^{-31^2/2}$. Adding independent dummy variables for unobserved atoms and applying exponential Markov yields, for all sufficiently large $n$,
\[
 \|b-b_A\|_n^2\le2^{-80}a^2
\]
except on probability at most $\exp(-2^{-83}n)$.

\emph{The posterior reference.}
Integrate the affine Gaussian block and use the same matrix $B$ as in Theorem~\ref{thm:all-scale-affine-likelihood}. Apply Lemma~\ref{lem:clipped-affine-prox-comparison}. The bounded function $u_A$, extended by zero off observed atoms, admits the uniform neural comparison with amplitude cap $32a$. Its entropy is bounded by
\[
 B_0+4\|u_A\|_\infty+e_2,
\]
where $B_0=\kappa(\nu)(H-1)$ is the depth baseline and $e_2$ has exactly the previous form, with a larger absolute coefficient on $a^2$. Its expected empirical loss differs from the loss of $u_A$ by at most $70a\varepsilon+\varepsilon^2$ at RMS accuracy $\varepsilon$.

Use this bounded comparison in the baseline-plus-amplitude proof, but retain the original proximal minimizer $u_b$ as the reference. The preceding lemma adds at most $2^{-80}a^2$ to its objective comparison. Explicitly, the low-temperature approximation error in the $B$ norm is at most
\[
 u+\sqrt{6au+s+2^{-80}a^2+\zeta+
             \lambda\{32a(4-c_\nu)+e_0+e_2\}}.
\]
Here $u$ is the fixed high-output threshold, $c_\nu=2\nu t$, and $\zeta=70a\varepsilon+\varepsilon^2$. This follows from the same two-comparison cancellation: the high-output quadratic gap is taken around $u_b$, while the comparison objective at $u_A$ is within the stated clipping error of that minimum.

All parameter estimates for the single cutoff
$\lambda_0=\Lambda a^2/(B_0+a)$ remain valid; only fixed numerical coefficients have increased. At high temperatures the zero comparison is at most $9a^2+s$, and the $B$-norm second-moment estimate is at most $40a^2$ after enlarging the width threshold. The indicator and amplitude branches therefore work exactly as before. Choose the constants and comparison tolerance so that the full posterior mean is within $2^{-31}a$ in training norm of a convex combination of the pure affine ridge reference and
\[
 f_b=(I-B)b+Bu_b.
\]
The original proximal minimizer may have a large cap, but its capped-affine representation and sector inequality remain valid for this finite Gaussian sample. No uniform cap on $u_b$ is required.

\emph{Uniform reference error and the common witness.}
The small-cap branch of Lemma~\ref{lem:capped-affine-error} is unchanged because the input Gram matrix and the target are unchanged. For its large-cap branch, use singleton Gaussian noises whose sign matches $\operatorname{sgn}(h-g)$ and whose magnitude is at least $a/2$. Each fixed sign pattern has this property at each singleton with probability at least one quarter. A Chernoff bound and the same hyperplane-pattern union bound give at least $n/24$ such singletons simultaneously for every affine $h$. At cap $t\ge a/1024$, each gives error at least $a/1024$. Consequently every capped-affine reference has squared empirical error at least $2^{-26}a^2$.
The Gaussian alignment obeys
\[
 |\langle b-g(S),g(S)\rangle_n|\le2^{-29}a^2
\]
with probability tending to one, since its conditional variance is at most $a^4/n$. The sector identity now yields the common reference witness at least $2^{-27}a^2$. The posterior approximation reduces it by at most $3\cdot2^{-31}a^2$, leaving
\[
 \langle b-2g(S),f-g(S)\rangle_n\ge2^{-28}a^2.
\]
This witness is common to every covariance and other learner parameter, so it survives arbitrary mixtures.

Finally, conditional on the input counts, the population witness
$h_S(x)=\sum_{i:X_i=x}(Y_i-2g(x))$ is
$a\sqrt{k_x}Z_x-k_xg(x)$. On the occupancy event,
\[
 \frac1n\sum_x k_xZ_x^2\le2
\]
with probability tending to one: its conditional mean is one and its variance is $2\sum_xk_x^2/n^2\le8/n$. Thus
\[
 \|h_S\|_{L^2(P_X)}^2
 \le\frac{2a^2}{n}\sum_xk_xZ_x^2
       +\frac2n\sum_xk_x^2g(x)^2\le12a^2<16a^2.
\]
The empirical-to-population witness identity is unchanged. Cauchy--Schwarz gives excess at least $2^{-60}a^2$, leaving slack for the displayed $2^{-64}a^2$. Increase the depth threshold so that the total sample failure is below one tenth. The actual learner always conditions on the original Gaussian labels; the clipping is used solely to construct its variational comparison.
\end{proof}



\section{Interpretation and scope}
At canonical mean field, a small hidden variance decays quadratically: $q-F(q)=\Theta(q^2)$. A relative covariance spike of order $q$ reverses this drift at entropy cost $\Theta(q^2)$ per layer. The sum along a reversed decay trajectory is finite and, when the endpoint variance is of order $a$, proportional to $a$. The finite-set construction first aligns this variance with a chosen signal; the cost of that alignment is quadratic in the logarithm of the number of encoded values and inversely proportional to depth.

The resulting function cost depends primarily on maximum amplitude on the specified finite domain. The learning theorem uses a separate statistical argument: the posterior approaches the empirical least-squares predictor with a common output cap, whose predictions retain the signs of noisy singleton labels. Those signs supply a common population-error witness at every temperature. The upper and lower prior estimates are both essential; a one-input tail estimate alone would not establish the learning obstruction.

Every learner prior in this paper is an independent Gaussian law in the family stated in Theorem~\ref{thm:all-scale-learning}; the critical member is \eqref{eq:canonical-prior}. Dependence enters only through the explicitly constructed comparison distributions. The proofs are finite-dimensional: they do not replace the network by a Gaussian process, take width to infinity, or restrict the prior's support. Each base learner uses one fixed dense architecture; the lower bound also covers the explicitly stated mixtures of their predictions. They have undergone separate AI-agent mathematical review but have not been formalized in Lean.
\appendix
\section{A limitation of one-direction covariance transport}
\label{sec:spike-method-barrier}
The square-root-depth scale in the covariance construction cannot be improved merely by optimizing its time-dependent spike factors. The obstacle is the entropy needed to amplify an exponentially small coherent seed.

\begin{proposition}[A quadratic action cost for coherence amplification]
\label{prop:coherence-action-barrier}
Consider the scalar covariance-spike recursion with arbitrary positive factors $R_\ell^2$ and valid states $0<r_\ell^2\le q_\ell$. Put $\theta_\ell=r_\ell^2/q_\ell$ and let $L$ be the number of transitions. Then
\[
 \sum_{\ell=1}^L
       \{R_\ell^2-1-\log R_\ell^2\}
 \ge\frac1{2L}
       \left(\log\frac{\theta_L}{\theta_0}\right)_+^2.
\]
In particular, raising coherence from $e^{-cD}$ to a fixed positive constant requires action $\Omega(D^2/L)$, for any deterministic time schedule. Allowing variance-reducing spikes $R_\ell<1$ does not evade the bound.
\end{proposition}
\begin{proof}
Write $t=R^2$, $v=q+(t-1)r^2$. Gaussian Stein and Cauchy--Schwarz give $v\chi(v)^2\le F(v)$. Thus
\[
 \frac{\theta'}\theta
 =\frac{t}{1+(t-1)\theta}\frac{v\chi(v)^2}{F(v)}
 \le\begin{cases}t,&t\ge1,\\1,&0<t<1.\end{cases}
\]
Consequently $\log(\theta_L/\theta_0)\le\sum_\ell(\log R_\ell^2)_+$.
For $u\ge0$, $e^u-1-u\ge u^2/2$. The action is therefore at least
$\tfrac12\sum_\ell((\log R_\ell^2)_+)^2$, and Cauchy--Schwarz proves the claim.
\end{proof}

\begin{lemma}[Some finite labels have tiny correlation with every first feature]
\label{lem:first-feature-discrepancy}
For any $m$ points in $\R^d$ with $m\ge d+2$, there is a sign function $f$ on those points satisfying
\[
 \sup_{a\in\R^{d+1}}
 \left|\frac1m\sum_x f(x)\tanh(a\cdot(1,x))\right|
 \le C\sqrt{\frac{(d+1)\log(2m)}m}.
\]
The same $f$ has this bound simultaneously for every finite collection of first-layer rows, regardless of their norms or the number of neurons.
\end{lemma}
\begin{proof}
The number of affine-halfspace sign patterns on the $m$ points is at most
$2\sum_{j=0}^{d}\binom{m-1}{j}\le(2em)^{d+1}$. This follows by the usual hyperplane-arrangement induction, and remains an upper bound for nongeneric point configurations. Choose independent fair signs $f(x)$. Hoeffding and a union bound show that with positive probability their normalized correlation with every such pattern has magnitude at most the displayed bound, for a sufficiently large absolute $C$.
Finally
\[
 \tanh t=\frac12\int_{\R}\operatorname{sign}(t-b)
                         \operatorname{sech}^2b\,db
\]
expresses each tanh feature as an average of affine-halfspace sign features. This transfers the correlation bound to all tanh rows at once.
\end{proof}

\begin{corollary}[Worst-case seeds impose the same barrier]
\label{cor:worst-seed-barrier}
On the unit-sphere sets $\mathcal U_D$ with $m=4\cdot2^D$, there are sign labels for which the following holds. Suppose a one-direction covariance-transport construction starts with arbitrary first-layer tanh features whose empirical variance at each input is at least $c/(D+1)$, chooses any unit seed direction, and then uses the scalar spike recursion with a common schedule. If its final scalar output is $a f(x)$ at every input, where $a>0$ is fixed and the final readout factor is bounded by a fixed $B_0$, its covariance action over $H$ transitions is at least $c' D^2/H$ for all sufficiently large $D$.
This conclusion permits arbitrary pointwise prewarping of seed magnitudes. It is a limitation of this transport construction, not a lower bound for every possible comparison measure or every Bayesian learner.
\end{corollary}
\begin{proof}
Choose $f$ from Lemma~\ref{lem:first-feature-discrepancy}. If $h_i(x)$ are arbitrary first features and $e$ is any unit direction, the seed $r_1(x)=N^{-1/2}\sum_i e_i h_i(x)$ obeys
\[
 \left|\frac1m\sum_x f(x)r_1(x)\right|
 \le\left\{\frac1N\sum_i
       \left(\frac1m\sum_x f(x)h_i(x)\right)^2\right\}^{1/2}
 \le C D/\sqrt m.
\]
Every scalar update multiplies $r$ by a positive number. After absorbing any common sign into the readout, the exact target therefore forces $r_1(x)$ to have sign $f(x)$. The left side is then the average of $|r_1(x)|$. At least one input has $|r_1(x)|\le CD/\sqrt m$, and hence initial coherence at most $C'D^3/m$ by the variance lower bound. At that same input the final correlation has magnitude at least $a/B_0$. Since the final variance is at most one, its final coherence is at least $a^2/B_0^2$. Proposition~\ref{prop:coherence-action-barrier} gives action at least
\[
 \frac1{2H}\left[\log\frac{c''m}{D^3}\right]^2
 =\Omega(D^2/H).
\]
No assumption that seed magnitudes are constant was used.
\end{proof}

\end{document}
